% concatenated from main.tex, macros.tex
\documentclass{article}
\usepackage{verbatim}
%\usepackage[utf8]{inputenc}
\newcommand{\R}{\mathbb{R}}
\newcommand{\Z}{\mathbb{Z}}
\newcommand{\C}{\mathbb{C}}
\usepackage{array}
\usepackage{natbib}
\usepackage{amsfonts}
\usepackage{bbm}
\usepackage{booktabs}
\usepackage{amsmath, mathrsfs}
\usepackage{mathtools}
\usepackage{enumitem}
\usepackage{amsthm}
\usepackage{xcolor}
\usepackage{graphicx} % Required for inserting images
\graphicspath{ {./figures/} }
\usepackage{caption}
\usepackage{subcaption}
\newtheorem{theorem}{Theorem}[section]
\newtheorem{corollary}{Corollary}[theorem]
\newtheorem{lemma}[theorem]{Lemma}
\newtheorem{example}[theorem]{Example}
\newtheorem{remark}[theorem]{Remark}
\newtheorem{proposition}[theorem]{Proposition}

\newtheorem{conj}[theorem]{Conjecture}
\DeclareMathOperator{\diag}{diag}

\usepackage{hyperref}

\newcolumntype{C}[1]{>{\centering\let\newline\\\arraybackslash\hspace{0pt}}p{#1}}
\newcommand{\unit}[1]{\ensuremath{\mathop{}\!\mathrm{#1}}}
\renewcommand{\Re}{\operatorname{Re}}
\renewcommand{\Im}{\operatorname{Im}}
\renewcommand{\angle}{\operatorname{angle}}
\newcommand{\diff}{\mathop{}\!\mathrm{d}}
\newcommand{\N}{\ensuremath{\mathcal{N}}}
\renewcommand{\mod}[1]{\ensuremath{\mathop{}[#1]}}

\newcommand{\floor}[1]{\ensuremath{\left\lfloor#1\right\rfloor}}

\DeclareMathOperator{\sinc}{sinc}
\DeclareMathOperator{\Spec}{Spec}
\DeclareMathOperator{\DFT}{DFT}

\newcommand{\set}[1]{\ensuremath{\left\{#1\right\}}}
\newcommand{\seq}[1]{\ensuremath{\left(#1\right)}}
\newcommand{\norm}[1]{\ensuremath{\left\lVert#1\right\rVert}}
\newcommand{\abs}[1]{\ensuremath{\left\lvert#1\right\rvert}}
\newcommand{\ip}[1]{\ensuremath{\left\langle#1\right\rangle}}
\newcommand{\fourier}{\ensuremath{\mathcal{F}}}
\newcommand{\frameset}{\ensuremath{\mathcal{F}}}
\newcommand{\gabor}{\ensuremath{\mathcal{G}}}
\newcommand{\defeq}{\vcentcolon=}



\newenvironment{notice}{%
\begin{center}
\color{red}
}{%
\end{center}
}
\newcommand{\textnotice}[1]{\textcolor{red}{(#1)}}

\newcommand{\cG}{\mathcal{G}}
\newcommand{\cT}{\mathcal{T}}

\usepackage{mathtools}
\usepackage{tikz}
\usetikzlibrary{matrix}

\usepackage[margin=1in]{geometry}




\title{On the structure of the Gram matrix for Gabor systems generated by B-splines.}

\date{\today}

%\begin{comment}
\author{
Martin Buck\thanks{Department of Mathematics, Tufts University, Medford MA 02131, USA. Email: martin.buck@tufts.edu}
\and
Christina Frederick\thanks{Department of Mathematical Sciences, New Jersey Institute of Technology, Newark NJ 07102, USA. Email: christina.frederick@njit.edu}
\and
Kasso A.~Okoudjou\thanks{Department of Mathematics, Tufts University, Medford MA 02131, USA. Email: kasso.okoudjou@tufts.edu}
\and
Alexander Stangl\thanks{Department of Mathematics, New Jersey Institute of Technology, Newark NJ 07102, USA. Email: ajs282@njit.edu}
}

%\end{comment} 

\begin{document}

\maketitle




\abstract{We consider the Gabor system $\mathcal{G}(g,a\mathbb{Z}\times b\mathbb{Z})$ generated by a continuous, compactly supported function $g$ over the time-frequency lattice generated by the parameters $a$ and $b$. We show that, under an appropriate ordering of the Gabor elements, certain submatrices of the Gram matrix of $\mathcal{G}(g,a\mathbb{Z}\times b\mathbb{Z})$ exhibit a block-Toeplitz structure. This structural property enables us to derive spectral results for finite sub-blocks of the Gram matrix  by appealing to the spectral theory of Toeplitz matrices. In particular, we apply our results to the Gram matrix of Gabor systems generated by the $N$th-order B-spline.


}



\section{Introduction}
\subsection{Background on the frame set problem}
Gabor analysis originated from two distinct perspectives: the operator-theoretic framework of J.~von Neumann \cite{MR66944} in 1932 and the communication-theoretic approach of D.~Gabor \cite{gabo1946} in 1946. Both authors independently investigated the completeness of the system 
\begin{equation}\label{gabgauss}
	\mathcal{G}(g,\mathbb{Z}^2) = \{g_{(n,k)}(\cdot) = g(\cdot - n) e^{2\pi ik \cdot} : (n,k) \in \mathbb{Z}^2\},
\end{equation}
utilizing the Gaussian $g(x)=e^{-\pi x^2}$ for its optimal localization in the time-frequency plane. This dual legacy highlights the fundamental nature of the integer lattice $\mathbb{Z}^2$ in tiling phase space (see \cite{heil2007history} for a detailed history).


More generally, the \textit{Gabor system} generated by a non-zero window function
\(g \in L^{2}(\R)\) along a (countable) discrete set of time-frequency shifts
\(\Gamma \subset \R^2\),
\begin{equation}\label{eq:Gabor_glambda}
	\mathcal{G}(g, \Gamma) = \{g_\gamma:=g(\cdot-\gamma_1) 
		e^{2\pi i\gamma_2\cdot}:\, \gamma=(\gamma_1,\gamma_2) \in \Gamma\},
\end{equation} is said to be a frame provided that there exist constants $0<\alpha\leq \beta < \infty$  such that for each \(f \in
L^2(\R)\)
\begin{equation}
	\alpha\norm{f}^2 \leq \sum_{\gamma \in \Gamma}
		\abs{\ip{f, g_{\gamma}}}^2 \leq \beta\norm{f}^2.
\end{equation}






    

It follows that there exists an associated set of functions 
\(\set{h_{\gamma}}_{\gamma \in \Gamma}\) which is also a Gabor frame for \(L^2(\R)\) for which the following reconstruction formulas hold for each \(f \in L^2(\R)\):
\begin{equation}
	f = \sum_{\gamma \in \Gamma} \ip{f, h_{\gamma}} g_{\gamma}
	  = \sum_{\gamma \in \Gamma} \ip{f, g_{\gamma}} h_{\gamma}.
\end{equation}
 When $\Gamma =a \Z \times b \Z$ is a separable lattice, it can be shown that the set of functions \(\set{h_{\gamma}}_{\gamma \in \Gamma}\) is a Gabor frame generated by a window function $h \in L^2(\R)$, and called a \textit{dual window} to $g$. In this case,  \(\set{h_{\gamma}}_{\gamma \in \Gamma}\) is a \textit{dual Gabor frame} to \(\set{g_{\gamma}}_{\gamma \in \Gamma} \)

 The full understanding of all window functions $g$, and all positive parameters $a, b$ such that $\mathcal{G}(g, a\mathbb{Z}\times b\mathbb{Z})$ is a Gabor frame for $L^2(\R)$, is still incomplete. Nonetheless, the last few decades of research have resulted in an impressive body of work on Gabor systems. For example, there exists a Nyquist-type condition dating back to the work of Daubechies~\cite{daubechies1990wavelet}, in which it is proved that when $ab >1$ is rational then $\mathcal{G}(g, a\mathbb{Z}\times b\mathbb{Z})$ is incomplete, for the irrational case we refer to \cite{MR1096368}. Furthermore, if $\mathcal{G}(g, a\mathbb{Z}\times b\mathbb{Z})$ is a Gabor frame for $L^2(\R)$ then $ab \leq 1$. The critical case $ab =1$ corresponds to the case when it is possible for $\mathcal{G}(g, a\mathbb{Z}\times b\mathbb{Z})$ to be an orthonormal or a Riesz basis. In such a case, the window function is necessarily poorly localized in the time-frequency plane as dictated by the Balian-Low Theorem \cite{MR1350699}. 
We refer to \cite{heil2007history} for an history of the density condition in Gabor analysis. We also refer to some of the foundational work of I.~Daubechies, who in the setting of  coherent state frames in the Bargmann-Fock conjectured that a Gabor system $G(g,a,b)$ on the rectangular lattice $a\Z\times b\Z$ is a frame if and only if $ab<1$ for the Gaussian, \cite{daubechies1988frames}. 


Still, it is in general unknown how to characterize the so-called
\textit{full frame set} 
\(\frameset(g)\) of a given window \(g \in L^2(\R)\)
\begin{equation}
	\frameset(g) = \set{\;\Gamma \subset \R^2 : \mathcal{G}(g, \Gamma) =
		\set{g_{(\gamma_1, \gamma_2)}}_{(\gamma_1,\gamma_2) \in \Gamma}
		\text{ is a frame over } L^2(\R)\;}.
\end{equation}


 






%\cite{daubechies1988frames}, Daubechies contributed greatly to the
%foundations of the theory of frames and their applications to the study
%of Gabor systems \cite{daubechies1986painless}. Notable results of these
%studies include the discovery of a density criterion required for
%completeness of a Gabor system, as well finding a region of the frame
%set for Gabor systems of continuous, compactly supported window
%functions for which for which a frame is obtained, the so-called
%``painless expansion''.}



% \cite{daubechies1986painless} This is painless nonorthogonal expansions, frame operator is multiple of identity

% \cite{daubechies1991simple} wilson bases with jaffard

% \cite{daubechies1994gabor} Links the analysis and synthesis operators in 

% Rewriting to emphasize shift in focus to wavelets following 1990 paper, emphasizing that we are building off of earlier work


% {\it Ingrid Daubechies' seminal contributions have profoundly shaped modern harmonic analysis, particularly by providing the rigorous mathematical framework for time-frequency representations crucial to signal processing. Her work with Grossmann and Meyer in \cite{daubechies1986painless}, was pivotal in formalizing the concept of frames and shifted the focus beyond traditional orthogonal bases, opening the door for practical systems like Gabor families. Later, in \cite{daubechies1988frames} with Grossmann, she further explored ``coherent state frames'' in the Bargmann space, leading to the significant conjecture that a  Gabor system generated by the Gaussian function forms a frame if and only if the density condition $ab<1$ is satisfied. Her landmark paper \cite{daubechies1990wavelet} then laid extensive groundwork for the theory of Gabor systems for continuous functions sampled on lattices. In it, she addressed density considerations, demonstrating in several important cases, that Gabor systems with $ab>1$ cannot span $L^2(\mathbb{R})$, thereby linking the Gabor lattice parameters directly to the completeness of the system. While the full characterization of Gabor frames generated by the Gaussian function was later established by others, her work provided the essential initial steps. 

% These investigations define the conditions ($ab<1$) under which the associated frame operator (whose properties are reflected in the Gram matrix) becomes invertible, guaranteeing stable signal reconstruction, and delineate when it becomes singular ($ab>1$). Furthermore, her work with Jaffard and Journe on Wilson bases in \cite{daubechies1991simple} demonstrated the construction of simple orthonormal bases with excellent time-frequency localization, providing a valuable contrast to the typically redundant nature of Gabor frames while still being rooted in the broader quest for efficient signal representations.




Restricting
\(\Gamma\) to the lattice \(\Lambda = a\Z \times b\Z\) for positive
constants \(a,b\), we arrive at what is referred to as the
\textit{frame set problem}. Insights offered by this variant of the
full frame set problem are hoped to enable further investigation of the
more general problem. 



Celebrated results in the early 90s characterized the frame set for the
Gaussian window as the open region $\{(a, b) \in \R^2_+: \, 0<ab<1\}$;  see \cite{seip1992density, seip1992density_bis, Lyubarskii92}. 
More recently, a similar result was obtained for classes of functions, including totally positive functions,  \cite{Grosto, Kloosto, grochenig2018sampling, grochenig2023totally}. For an overview of the state of this fascinating problem we refer to \cite{Olehon, grochenig2014mystery}. For some historical perspective, we refer to the work of Daubechies, Grossmann and Meyer, \cite{daubechies1986painless},   on painless non-orthogonal expansions from which one concludes that the rectangle $(0, 2)\times (0, 1/2)$ belongs to the frame set for many window functions $g$. 


For $g\in L^2(\R), g\neq 0$, and $f\in L^2(\R)$, let $V_gf$ be the function defined on $\R^2$ by $$V_{g}f(x, \omega)=\int_{\R}f(t)\overline{g(t-x)}e^{-2\pi i t\omega}\, dt.$$ The modulation space $M^{1}(\R)$, also known as the Feichtinger algebra $S_0(\R)$, is the dense Banach subspace of $L^2(\R)$ consisting of all functions $g$ such that $V_gg\in L^1(\R^2)$ \cite{BenOko20, Groc2001}. It was proved in \cite{FeiKai04} that if the window function $g$ is in the modulation space $ M^1(\R)$  then its frameset $\mathcal{F}(g)$ is an open subset of $\{(a, b) \in \R^2_+: \, 0<ab<1\}$. 



Examples of functions in $M^1(\mathbb{R})$ are the B-splines of order $N=2$ or higher, which are  even, compactly supported functions defined
recursively as
\begin{equation}\label{eq:sn}
	s_1(x) = \chi_{[-1/2, 1/2]}(x), \qquad s_N(x) = s_1(x) * s_{N-1}(x).
\end{equation}
They are of special interest to study for a number of reasons:
B-spline wavelets, which have a construction similar to Gabor
systems, have been enormously successful in finite element methods
\cite{xiang2006construction, jiawei2008new, xiang2006identification, wei2021b}. Gabor systems
corresponding to B-spline window functions have also shown promise
in acoustic scattering problems \cite{Kreuzer2019}.

The frame set for the first order B-spline $s_1$ was completely
characterized in \cite{dai2016abc}, and work on higher order
B-spline has gradually accumulated in the past decade (see
\cite{AtiKouOko1} \cite{atindehou2023frame} \cite{Olehon} \cite{Olehon1}
\cite{Ole4} \cite{ghosh2025gabor} \cite{ghosh2023obstructions} 
\cite{Grojan} \cite{Lemniel} \cite{Kloosto} for details). As noted above,
 \(\frameset(s_2)\) is an open set in $\{(a, b) \in \R^2_+: \, 0<ab<1\}$ \cite{FeiKai04},
but a full characterization of this set remains an open question. 

%In an
%earlier work \cite{AtiKouOko1}, the authors 

The authors of \cite{AtiKouOko1} introduced a framework for
determining the frame sets of compactly supported functions, including
the B-spline of order \(N \geq 2\), that unified many known results
on the frame sets of B-spline. Similarly, the authors of  \cite{atindehou2023frame} 
%the authors 
used a similar linear algebra based approach to shed new light
on the set \(\frameset(s_2)\) and enlarged the known frame set region by
deriving sufficient conditions for invertibility of the linear system
relating a frame with its dual. The key idea is based on the following result (see \cite[Theorem
2.10]{Ole1}): for \(g,h \in L^2(\R)\) and \(a,b > 0\), the Bessel
sequences \(\mathcal{G}(g, a\Z \times b\Z)\) and \(\mathcal{G}(h, a\Z \times
b\Z)\) are dual frames for \(L^2(\R)\) if and only if
\begin{equation}\label{eq:dualgeneral}
	\sum_{k \in \Z} \overline{g(x-\ell/b+ka)} h(x+ka) =
		b\delta_{\ell,0}, \forall \ell \in \Z \mbox{ and a.e. } x \in
		\left[-\frac{a}{2},\frac{a}{2}\right].
\end{equation}
The authors of \cite{atindehou2023frame} showed the existence of a
bounded compactly supported function \(h \in L^2(\R)\) that solves
\eqref{eq:dualgeneral} when \((a,b)\) belongs to a family of sets
parameterized by \(m\). To do this, the authors simplified and rewrote
\eqref{eq:dualgeneral} as a matrix-vector equation using the compact
support of \(s_2\) and derived sufficient conditions for invertibility
of the system.






\subsection{Contributions}

This paper aims to investigate the frameset of continuous, compactly supported functions. To this end, we initiate an analysis of the Gram matrix of the corresponding Gabor systems.
We build upon the structural investigations introduced in \cite{buckspectral}, and provide a spectral and structural analysis of the Gram matrices of Gabor systems generated by a class of continuous compactly supported windows, with a specific focus on B-splines of order $N$. In particular, we identify a unique structural hierarchy within these Gram matrices; we demonstrate in Theorem \ref{thm:offdiagTH} that certain a finite sub-Gramian  is a block-Toeplitz matrix where each constituent block admits a Hadamard factorization into a real-valued Toeplitz matrix and a rank-one Hankel matrix. Secondly, we derive an explicit analytic formulation for the Laurent symbols $t^{[\ell]}(x)$ associated with these blocks, proving in Theorem \ref{thm:tl_sinc} that they are governed by periodic sums of $\text{sinc}^{2N}$ kernels. Using this formulation, we establish an $O(\ell^{-N})$ spectral decay rate for the off-diagonal blocks, which provides the first formal mathematical justification for the block-diagonal dominance observed in the numerical analysis of spline-based Gabor systems. Finally, we leverage the theory of asymptotically equivalent circulant matrices to provide a pathway for estimating spectral bounds, specifically demonstrating how rational lattice parameters $a \in \mathbb{Q}$ induce the decay of the lower frame bounds in truncated systems.

The remainder of this paper is organized as follows. Section~\ref{sec:notations} sets the notations use throughout the paper. Section \ref{sec:gram_block} establishes the internal block structure of the Gram matrix and provides the proof for the Toeplitz-dot-Hankel factorization. Section \ref{sec:asymptotic} presents the spectral analysis of the block operators $\mathscr{G}^{[\ell]}$ via their associated Laurent symbols. We also introduce a sequence of asymptotically equivalent circulant matrices used for finite-section eigenvalue estimation. In addition, we discuss the implications of these spectral results on the frame set problem, specifically regarding the behavior of frame bounds on rational and irrational lattices. Finally, Section \ref{sec:conclusion} summarizes our findings and proposes future directions.









\section{Notation and preliminaries}\label{sec:notations}


Given a discrete collection of functions 
\(\set{g_\lambda}_{\lambda \in \Lambda}\) of \(L^2(\R)\), we define the following operators: 
\begin{itemize}
	\item The \textit{synthesis operator} \(\mathscr{T}\) is a mapping from
	      \(\C^{\Lambda}\) to \(L^2(\R)\) defined as
		  \begin{equation}
			  \mathscr{T}\seq{c_{\lambda}}_{\lambda \in \Lambda} = \sum_{\lambda
			      \in \Lambda} c_{\lambda} g_{\lambda}.
		  \end{equation}
	\item Its (formal) adjoint is the \textit{analysis operator} \(\mathscr{T}^*\) is given by
		  \begin{equation}
		      \mathscr{T}^*f = \seq{\ip{f, g_{\lambda}}}_{\lambda \in \Lambda}.
		  \end{equation}
	\item The \textit{frame operator} \(\mathscr{S}\) defined on
	      \(L^2(\R)\) is the composition of the synthesis and
		  analysis operators and given by  \(\mathscr{T}\mathscr{T}^*\)
		  \begin{equation}
		      \mathscr{S}f = \sum_{\lambda \in \Lambda} \ip{f, g_{\lambda}}
			      g_{\lambda}.
		  \end{equation}
	\item The \textit{Gram operator} \(\mathscr{G}\) is defined on 
	      \(\C^{\Lambda}\) as the composition of the analysis and
		  synthesis operators  \(\mathscr{T}^*\mathscr{T}\) and given  by
		  \begin{equation}
			  \mathscr{G}\seq{c_{\lambda}}_{\lambda \in \Lambda} =
			      \seq{\sum_{\lambda' \in \Lambda} c_{\lambda'}
				  \ip{g_{\lambda'}, g_{\lambda}}}_{\lambda \in \Lambda}.
		  \end{equation}
		  The matrix corresponding to \(\mathscr{G}\) is
		  referred to as the \textit{Gram matrix}. 
\end{itemize}
The frame operator and the Gram operator are related in that the two
share a spectrum, with the possible exception of zero, which may be
present in the spectrum of the Gram operator but not that of the frame
operator. For this reason, while the frame operator may be more readily
amenable to a finite sections method the presence of zero in the
spectrum of the Gram operator means that the spectra of the finite sections of the Gram
operator do not necessarily converge to the non-zero portion of the spectrum of the Gram operator (see \cite{christensen2005finite}).
This manifests as \textit{pollution} filling the spectral gap between
zero and \(\norm{S^{-1}}^{-1}\) when taking finite sections of \(\mathscr{G}\).







For the rest of the paper the discrete systems we consider are Gabor systems $\mathcal{G}(g, \Lambda)$ generated by a window function $g\in L^2(\R)$ and where $\Lambda$ is the lattice $a\Z \times b\Z$. We summarize in Table~\ref{tab:notation} below the notations to be used in the sequel. 

\begin{center}
	\renewcommand{\arraystretch}{1.3}
	\begin{tabular}{cl}
		\toprule
		\textbf{Symbol} & \textbf{Description} \\
		\midrule
		\(\mathcal{G}(g,S)\) & Gabor system \(\set{T_y M_{\omega} g :
			(y,\omega) \in S}\) with \(T_y f(x) = f(x-y)\), \(M_{\omega}f(x)
			 = e^{2\pi i\omega x} f(x)\) \\
		\(\Lambda\) & Infinite lattice \(a\Z \times b\Z\)\\
		\(\mathscr{G}\) & Infinite Gram matrix/operator over
			\(\Lambda\) \\
		\hline
		\(\Lambda_n\) & Finite shifted lattice \(\set{(ka,jb) :
			j,k = -\tfrac{n-1}{2}, \hdots, \tfrac{n-1}{2}}\) \\
		\(g_{j,k}\) & Atom \(g_{j,k}(x) \defeq T_{ka}M_{jb}g(x) =
			g(x-ka)e^{2\pi ijb x}\) \\
		\(G_n\) & Truncated \(n^2 \times n^2\) Gram matrix over
			\(\Lambda_n\) \\
		\(G_n^{[\ell]}\) & \(n \times n\) block at modulation difference
			\(\ell\) (\(\ell = j-j'\)), \(\abs{\ell} < n\) \\
		& Entries \((G_n^{[\ell]})_{k,k'} \defeq
			\ip{g_{j,k},g_{j-\ell,k'}}\), \(0 \leq k,k' \leq n-1\) \\
		\(\set{\lambda_m(C_n)}_{m=1}^n\) & Ordered eigenvalues of an
			\(n \times n\) matrix \(C_n\) \\
		\bottomrule
	\end{tabular}
	\vskip1em
	\par Table 1: Notation for time-frequency shifts, Gram matrices, and
	block structure.
	\label{tab:notation}
\end{center}

Fix \(a,b > 0\) with \(ab < 1\).  Let \(n \geq 3\) be an odd natural
number such that \(n > \frac{2}{a} + 1\), and consider  the
index set be \(\mathcal{I}_n \defeq \set{-\frac{n-1}{2}, \ldots,
\tfrac{n-1}{2}}\).

In the sequel we will focus on time-frequency shifts of the
$N$th order B-spline, \(g = s_N\), defined in \eqref{eq:sn}:
\begin{equation}
	g_{j,k}(x) = g(x - ka)e^{2\pi i jb x}, \qquad j,k \in \mathcal{I}_n
\end{equation}
in which the time-frequency shifts belong to a set of enumerated
elements,
\begin{equation*}
	\Lambda_n = \set{(ka, jb)}_{j,k \in \mathcal{I}_n}
\end{equation*}
The infinite Gram matrix corresponding to the Gabor system
\eqref{eq:Gabor_glambda} is defined entrywise by
\begin{equation}\label{eq:graminf_entry}
	(\mathscr{G})_{(j,k),(j',k')} \defeq \ip{g_{j,k}, g_{j',k'}} \qquad
		j, j', k, k' \in \Z.
\end{equation}

For each odd \(n \in \mathbb{N}\) with \(n \geq 3\) we introduce the truncated
system \(\set{g_{\lambda}}_{\lambda \in \Lambda_n}\). Then the sets
\(\Lambda_{2k+1}\) are nested:
\begin{equation*}
	\Lambda_3 \subset \Lambda_5 \subset \hdots, \qquad
		\bigcup_{k = 0}^\infty \Lambda_{2k+1} = \Lambda.
\end{equation*}

For the finite index set \(\mathcal{J}_n \defeq \mathcal{I}_n \times
\mathcal{I}_n\), let \(P_n : \ell^2(\Z^2) \to \ell^2(\mathcal{J}_n)\) denote
the coordinate projection onto \(\mathcal{J}_n\). Then the finite Gram
matrix \(G_n\) is the corresponding finite section
\begin{equation*}
	G_n = P_n \mathscr{G} P_n^*.
\end{equation*}
As \(n \to \infty\) the projections \(P_n\) converge strongly to the
identity, so \(G_n \to \mathscr{G}\) in the strong operator topology. The
entries of \(G_n\) are given by
\begin{equation}\label{eq:gram_entry}
	(G_n)_{(j,k), (j',k')} \defeq \ip{g_{j,k}, g_{j',k'}}, \qquad
		j, j', k, k' \in \mathcal{I}_n.
\end{equation}

We may view \(G_n\) as an \(n^2 \times n^2\) block matrix indexed by the
modulation indices \(j,j' \in \mathcal{I}_n\), where each block
\(G_n^{(j,j')} \in \C^{n \times n}\) corresponds to a fixed pair
\(j,j'\) and has entries
\begin{equation*}
	(G_n^{(j,j')})_{k,k'} = \ip{g_{j,k}, g_{j',k'}} \qquad
		k, k' \in \mathcal{I}_n.
\end{equation*}
This block structure will be exploited to describe the Toeplitz
properties and structural behavior of \(G_n\) and its infinite lattice
limit \(\mathscr{G}\).

Recall that a square \textit{Toeplitz matrix} \(A \in \C^{n \times n}\)
is a matrix whose entries obey the relation
\begin{equation*}
	A_{k,k'} = A_{k+1,k'+1} = a_{k-k'} \qquad
		k, k' = 0, 1, \ldots, n-1.
\end{equation*}

A \textit{Circulant matrix} is a Toeplitz matrix \(C \in \C^{n \times
n}\) having the form
\begin{equation*}
	C_{k,k'} = c_{k-k' \mod n} \qquad
		k, k' = 0, 1, \ldots, n-1.
\end{equation*}

A Hermitian matrix \(T\) has an \(n \times n\) \textit{block-Toeplitz}
structure with \(k \times k\) blocks if
\begin{equation*}
	T = \begin{pmatrix}
		A^{[0]} & A^{[1]} & \cdots & A^{[n-1]} \\[1ex]
		A^{[-1]} & A^{[0]} & \cdots & A^{[n-2]} \\[1ex]
		\vdots & \vdots & \ddots & \vdots \\[1ex]
		A^{[-(n-1)]} & A^{[-(n-2)]} & \hdots & A^{[0]}
	\end{pmatrix}
\end{equation*}
with
\begin{equation*}
	A^{[j]} \in \C^{k \times k}, \qquad A^{[j]} = (A^{[-j]})^*, \qquad
		j = -n+1, \hdots, n-1
\end{equation*}
If the matrices in the block-Toeplitz matrix are themselves Toeplitz,
then the matrix is \textit{Toeplitz-block-Toeplitz}.



\section{Toeplitz structures in the Gram matrix for $\mathcal{G}(s_N, a\mathbb{Z}\times b\mathbb{Z} )$}\label{sec:gram_block}



In this section we show first that the enumeration of $\Lambda_n$ can be done without loss of generality when considering the spectrum of ${G}_n$. This is because other orderings will be permutations of the rows and columns of ${G}_n$ that are simple similarity relations that leave the spectrum unchanged. Then, we show that $G_n$ has block structure with blocks that have an explicit representation as a Hadamard product of a Toeplitz and rank-one Hankel matrix. We then provide resulting properties of these matrices.


% We first justify the choice of centering the finite sub-lattice at the origin by showing that if one were to consider a square sub-lattice centered at another location in the time-frequency plane, the entries of ${G}_n$ change in a predictable way by a potential phase factor. The proof is a simple change-of-variables in the inner product seen in Equation \ref{eq: gram_entry}.


% {\color{red} justification for sampling in this way, not as strong as "it doesn't matter how you sample", how you sample the time-frequency lattice is not an exact science. instead of a block around the origin, if you have a block somewhere else, you can shift it back around the evidence, numerical evidence says it doesn't affect the spectrum. kasso thinks it should not. explicitly, square at a different location, shift the square at the center at the origin, there will be different gram matrix but is this a unitary transformation? same eigenvalues? alex: shifting in b doesn't change the gramian, shifting horizontally. prove that it is a unitary transformation so the spectrum isn't changed.}

% {\color{purple} The above discussion, formalized in this lemma, also says that $\Gamma^{[\ell]}$ and $G_n^{[\ell]}$ are different based on a phase factor}.
\begin{lemma} \label{lem:shift-invariance}
Let $s, t \in \mathbb{R}$ and define the shifted set:
\[
\widetilde{\Lambda_n} := \Lambda_n+(as, bt),
\]
and the corresponding Gabor atoms
\begin{equation}
\tilde{g}_{j,k}(x) = g(x - (ka+sa))e^{2\pi i (jb+tb) x}, \qquad j,k\in \mathcal{I}_n
\end{equation}

Then, for any $(j,k), (j',k') \in \Lambda_n$, the new Gram matrix $\tilde{G}_n$ has entries
\[
\langle \tilde{g}_{j,k}, \tilde{g}_{j',k'}\rangle = e^{2\pi i b as (j - j')} \langle g_{j,k}, g_{j',k'} \rangle.
\]
The shifting factor $s$ $(s=1/ab)$ can be chosen so that $\tilde{G}_n=G_n$ is unchanged. Furthermore, the spectra of $G_n$ and $\tilde{G}_n$ are identical.
\end{lemma}

\begin{proof}
The proof is a simple change-of-variables $u=x-as$ in the inner product seen in Equation \ref{eq:gram_entry}:
\begin{align*}
\langle \tilde{g}_{j,k}, \tilde{g}_{j',k'} \rangle
&= \int_{\mathbb{R}} g(x - ka - sa) e^{2\pi i (jb + tb)x} \cdot g(x - {k'}a - sa) e^{-2\pi i ({j'}b + tb)x} \, dx \\
&= \int_{\mathbb{R}} g(u - ka) e^{2\pi i jb (u + sa)} \cdot g(u - {k'}a) e^{-2\pi i {j'}b (u + sa)} \, du \\
&= e^{2\pi i b s a(j-j')} \int_{\mathbb{R}} g(u - ka) g(u - {k'}a) e^{2\pi i (jb - {j'}b) u} \, du \\
&= e^{2\pi i b s a(j-j')} \langle g_{j,k}, g_{j',k'} \rangle.
\end{align*}

Setting $U=\mathrm{diag}(e^{2\pi i a b s\, j})_{j=0}^{n-1}$,
\[
\tilde G_n=(U\otimes I_n)\,G_n\,(U\otimes I_n)^*,
\]
so $\tilde G_n$ is unitarily similar to $G_n$ and has the same spectrum.


\end{proof}


Having established that lattice shifts induce only phase modifications, we next analyze the internal block structure of $G_n$ itself. The below lemma shows that the finite Gram matrix $G_n$ is structured as an $n^2\times n^2$ block matrix whose $(j, j')$-th block is given by an $n\times n$ matrix $G_n^{[\ell]}$ with $\ell = j - j'$.

\begin{lemma} \label{lem: diag_blocks}
The matrix ${G}_n$ given by \eqref{eq:gram_entry} is block-Toeplitz with blocks that are banded for $n>\frac{2}{a}+1$. 
\end{lemma}



\begin{proof}

For $j, j', k, k'\in \mathcal{I}_n$, the $(j,j')-$th block in $G_n$
has its $(k,k')-$th entry given by
\begin{align}
\langle g_{j,k}, g_{j',k'} \rangle = \int_{\mathbb{R}} g(x - {ka}) g(x - {k'a}) e^{2\pi i ({jb} - {j'b})x} \, dx.
\end{align}
Due to the fixed modulation phase $e^{2 \pi i ({jb} - {j'b})x}$ appearing above, the blocks depend only on the modulation difference $\ell:=j-j'$, i.e. $G_n$ is
block–Toeplitz in $j$, and for $|\ell|\le n-1$ we write
    \begin{align*}
     (G_n^{[\ell]})_{k,k'}:=\langle g_{j,k}, g_{j-\ell,k'} \rangle = \int_{\mathbb{R}} g(x - {ka}) g(x - {k'a}) e^{2\pi i ({jb} - {(j-\ell)b})x} \, dx, 
    \end{align*}
This proves that $G_n$ is block-Toeplitz and has the form
\begin{align*}
        G_n &=
\begin{pmatrix}
G_n^{[0]} & G_n^{[1]} & G_n^{[2]} & \cdots & G_n^{[n-1]} \\
G_n^{[-1]} & G_n^{[0]} & G_n^{[1]} & \cdots & G_n^{[n-2]} \\
G_n^{[-2]} & G_n^{[-1]} & G_n^{[0]} & \cdots & G_n^{[n-3]} \\
\vdots & \vdots & \vdots & \ddots & \vdots \\
G_n^{[-(n-1)]} & G_n^{[-(n-2)]} & G_n^{[-(n-3)]} & \cdots & G_n^{[0]}
\end{pmatrix}.
\end{align*}
    
The compact support of $g=s_N$ on $[-\frac{N}{2},\frac{N}{2}]$ produces
\begin{align*}
    (G_n^{[\ell]})_{k,k'}&=
    \begin{cases}
e^{\pi i b l ({ka} + {k'a})}
\displaystyle\int_{{k'a} - 1}^{{ka} + 1} (x - {k'a} + 1)(-x + {ka} + 1) \cos(2\pi b\ell x) \, dx, & |{ka} - {k'a}| \leq N \\[2ex]
0, & \text{otherwise}
\end{cases}.
\end{align*}
Since the integral is zero when $|{k'a}-{ka}| > N$, or $|k-k'|>\frac{N}{a}$, the blocks ${G}_n^{[\ell]}$ are banded (recall that $n > \frac{N}{a}+1$).








% For the diagonal blocks, $k=0$, the integral for the case $|n-n'|<\frac{2}{a}$ is seen to be:
% \begin{equation}
%      G_n^{[0]}_{k,k'}= \frac{1}{6}({ka}-{k'a}+2)^3.
% \end{equation}
% Therefore, the entry $G_n^{[0]}_{k,k'}$ is real and a function of the quantity $n-n'$, hence Toeplitz. 

% Taking the magnitude or magnitude-squared eliminates the phase factor $e^{2\pi i bk {ka}}$ and the resulting entries in the blocks $(G_n^{[\ell]})_{k,k'}$ depend only on $l$ and $n-n'$; therefore $|G_n^{[\ell]}|^2$ is Toeplitz-block-Toeplitz.



\end{proof}





The next theorem describes the structure of the blocks given in Lemma \ref{lem: diag_blocks}


\begin{theorem} \label{thm:offdiagTH}
Let $G_n\in \mathbb{C}^{n^2\times n^2}$ be the truncated Gram matrix corresponding to a Gabor system generated by $g=s_N$ over $\Lambda_n$. Then each block $G_n^{[\ell]}\in \mathbb{C}^{n\times n}$, $|\ell|\leq n-1$ can be written as a Hadamard (entrywise) product:
\begin{align}
G_n^{[\ell]} = T_n^{[\ell]} \circ H_n^{[\ell]},\label{eqn:gram:toeplitz-dot-hankel}    
\end{align}

where $T_n^{[\ell]} \in \R^{n\times n}$ is a real-valued, symmetric Toeplitz matrix, and $H_n^{[\ell]}\in \C^{n\times n}$ is a rank-one Hankel matrix with unimodular entries.
\end{theorem}


\begin{proof}
For $k,k'\in \mathcal{I}_n$, we consider the $(k,k')-$entry of the block \(G_n^{[\ell]}\) corresponding to a fixed modulation difference $\ell$:
\[
(G_n^{[\ell]})_{k,k'}:=\langle g_{j,k}, g_{j-\ell,k'} \rangle = \int_{\mathbb{R}} g(x - {ka}) g(x - {k'a}) e^{2\pi i b \ell x} \, dx.
\]

After a symmetric change of variables \(x \to  x + \tfrac{1}{2}({ka} + {k'a})\), the entries are given by:
\begin{align*}
(G_n^{[\ell]})_{k,k'}
&= \int_{\mathbb{R}} g(x + \tfrac{1}{2}({ka} + {k'a}) - {ka})\, g(x + \tfrac{1}{2}({ka} + {k'a}) - {k'a})\, e^{2\pi i b \ell (x + \tfrac{1}{2}({ka} + {k'a}))} \, dx \\
%&= e^{4\pi i \overline{y} \Delta \omega} \int_{\mathbb{R}} g(u - \Delta y) g(u + \Delta y)\, e^{4\pi i u \Delta \omega} \, du \\
%&= e^{2\pi i bk \overline{y} } \int_{\mathbb{R}} g(x - \Delta y)\, g(x + \Delta y)\, \cos(4\pi \Delta \omega x) \, dx\\
&=e^{\pi i b\ell ({ka} + {k'a}) }\int_{\mathbb{R}} g\left(x - \tfrac{a}{2} (k-k')\right) g\left(x + \tfrac{a}{2} (k-k')\right)\, \cos(2\pi b\ell x) \, dx.
\end{align*}
Define the Toeplitz matrix $T_n^{[\ell]}$ with entries given by $(T_n^{[\ell]})_{k,k'}:= t^{[\ell]}_{\,k-k'}$, where, for $j=k-k'$
\begin{align}
    t^{[\ell]}_{\,j} &:= \int_{\mathbb{R}} g\left(x - \tfrac{a}{2} j\right) g\left(x + \tfrac{a}{2} j\right)\, \cos(2\pi b\ell x) \, dx, \label{eq:tjl}
\end{align}
and define $(H_n^{[\ell]})_{k,k'}:=e^{\pi i a b \ell (k+k')}$. Then $G_n^{[\ell]}=T_n^{[\ell]}\circ H_n^{[\ell]}$. Since $H_n^{[\ell]}=v^{[\ell]}(v^{[\ell]})^{\!\top}$ with
$v^{[\ell]}=\big(e^{\pi i a b \ell\,k}\big)_{k\in\mathcal{I}_n}$, it follows that $H_n^{[\ell]}$ is rank-one Hankel. Furthermore, since $g=s_N$ has support in $[-\frac{N}{2},\frac{N}{2}]$, we have $t^{[\ell]}_{\,j}=0$ for $|j|>\lfloor N/a\rfloor$, so
$T_n^{[\ell]}$ (hence $G_n^{[\ell]}$) is banded with bandwidth $\lfloor N/a\rfloor$.
\end{proof}

We can immediately relate the spectra of $G_n$ and that of $T_n^{[\ell]}$. 


\begin{corollary}\label{cor:specgntnl}
Each block $G_n^{[\ell]}$ defined in Theorem \ref{thm:offdiagTH} shares the same singular values as the corresponding Toeplitz matrix $T_n^{[\ell]}$. Consequently, the Schatten norms and condition numbers of $G_n^{[\ell]}$ and $T_n^{[\ell]}$ are identical.
\end{corollary}

\begin{proof}
Let \( H_n^{[\ell]} =  v^{[\ell]}(v^{[\ell]})^{\!\top} \) and $T_n^{[\ell]} \in \mathbb{R}^{n\times n}$ be the Hankel and Toeplitz matrices given in the proof of Theorem~\ref{thm:offdiagTH}, where \( v^{[\ell]} \in \mathbb{C}^n \) is a vector of unit-modulus entries. Define the diagonal matrix \( D_n^{[\ell]} := \mathrm{diag}(v^{[\ell]}) \). Then the block \( G_n^{[\ell]} \) can be further written as
\[
G_n^{[\ell]} = D_n^{[\ell]} T_n^{[\ell]} D_n^{[\ell]}.
\]



Since $D_n^{[\ell]}$ is unitary,
\[
G_n^{[\ell]} G_n^{[\ell]*} = D_n^{[\ell]} T_n^{[\ell]} (T_n^{[\ell]})^* (D_n^{[\ell]})^*,
\]
which has the same eigenvalues as $T_n^{[\ell]} (T_n^{[\ell]})^*$. Thus the singular values coincide.
\end{proof}

Matrices with the decomposition in Theorem \ref{thm:offdiagTH} were coined {\it Toeplitz-dot-Hankel} matrices in \cite{townsend2018fast} and used to develop fast algorithms for conversions between coefficients in different orthogonal polynomial expansions. The top two plots in Figure \ref{fig: gram_block} show the magnitude and phase of the matrix $G_n$, and the bottom two plots show block structures that are revealed by the decomposition \eqref{eqn:gram:toeplitz-dot-hankel}: the Toeplitz matrices $T_n^{[\ell]}$ (bottom left) and the phase of the Hankel matrices $H_n$ (bottom right).

\begin{figure}[h!]
    \centering
  \includegraphics[width=\textwidth]{Figures/G15.png}
     \caption{Top: The magnitude and phase of the finite Gram matrix $G_{15} \in \mathbb{C}^{115 \times 115}$ with blocks $G_{15}^{[\ell]} \in \mathbb{C}^{15 \times 15}$  Bottom: The magnitude of the matrix $T_{15}$ with Toeplitz blocks $T^{[\ell]}_{15}$ and the phase of the matrix $H_{15}$ with Hankel blocks $H_{15}^{[\ell]}$ in the decomposition \eqref{eqn:gram:toeplitz-dot-hankel}
     for $a=.25, b=1.5, n=15$.}
     \label{fig: gram_block}
\end{figure}




 By reordering indices according to modulation difference, $\mathscr{G}$ may also be expressed in block form with blocks $\mathscr{G}^{[\ell]}$ which allows us to extend Theorem~\ref{thm:offdiagTH} to $\mathscr{G}$.  For $\ell\in\mathbb{Z}$, the
$\ell$–block (grouping by modulation difference $j-j'=\ell$) is the
bi-infinite matrix on $\ell^2(\mathbb{Z})$ given by
\[
(\mathscr{G}^{[\ell]})_{k,k'} := \langle g_{j,k}, g_{j-\ell,k'}\rangle
= \int_{\mathbb{R}} g(x-ka)\,g(x-ka')\,e^{2\pi i b \ell x}\,dx.
\]





\begin{corollary}\label{cor:infBlocksTn}
Each bi-infinite block $\mathscr{G}^{[\ell]}$ admits a Hadamard (entrywise) factorization
\[
\mathscr{G}^{[\ell]} = T_{\infty}^{[\ell]} \circ H_{\infty}^{[\ell]},
\]
where $T_{\infty}^{[\ell]}$ is a banded, real-valued, symmetric Toeplitz operator on $\ell^2(\mathbb{Z})$
and $H_{\infty}^{[\ell]}$ is a complex-valued, rank-one Hankel operator with unimodular entries.
Moreover, letting $D_{\infty}^{[\ell]}=\mathrm{diag}(e^{\pi i a b \ell k})_{k\in\mathbb{Z}}$,
we have the diagonal unitary similarity
\[
\mathscr{G}^{[\ell]} = D_{\infty}^{[\ell]} \, T_{\infty}^{[\ell]} \, D_{\infty}^{[\ell]},
\]
so $\mathscr{G}^{[\ell]}$ and $T_{\infty}^{[\ell]}$ are unitarily equivalent and therefore share the same
spectra.
\end{corollary}



 
We recall that a matrix $A$ indexed by $\mathcal I\times\mathcal I$ is
\emph{per-Hermitian} if $A = J\,A^*\,J$, where $J$ is the reversal
(involution) $(Jx)_k := x_{-k}$ for symmetric indexing 
$k\in\mathcal I$ (and $(Jx)_k := x_{n-k+1}$ for $k\in\{1,\dots,n\}$).

\begin{proposition}\label{prop:perHermitian}
Let $g=s_N$ be the $N$th order B-spline. Then the finite Gram matrix $G_n$ and the infinite Gram operator $\mathscr{G}$
are per-Hermitian with respect to reversal in both indices:
\[
\overline{(G_n)_{(j,k),(j',k')}} \;=\; (G_n)_{(-j',-k'),(-j,-k)},
\qquad
\overline{(\mathscr{G})_{(j,k),(j',k')}} \;=\; (\mathscr{G})_{(-j',-k'),(-j,-k)}.
\]
Consequently, when $G_n$ is viewed as an $n\times n$ block matrix in the
modulation indices, each block $G_n^{[\ell]}$ is per-Hermitian:
\[
\overline{(G_n^{[\ell]})_{k,k'}} \;=\; (G_n^{[\ell]})_{-k',-k},
\qquad k,k'\in\mathcal I_n,
\]
and likewise for $\mathscr{G}^{[\ell]}$ on $\ell^2(\mathbb Z)$.
\end{proposition}


% \begin{proof}
%     Because a Gram matrix is Hermitian, we just need to show that $(G_n^{[\ell]})_{k,k'} = (G_n^{[\ell]})_{n-k+1, n-k'+1}$. Then, because the set $\{-\lfloor n/2 \rfloor, \ldots, \lfloor n/2 \rfloor\}$ is symmetric about the origin so that $|{ka}| = |y_{n-k+1}|$, we have
% \begin{align*}
%     (G_n^{[\ell]})_{n-k+1, n-k'+1} &= \int_{\mathbb{R}} g(x-y_{n-k+1})e^{2 \pi i {jb} x} g(x-y_{n-k'+1}) e^{-2 \pi i {j'b}x} dx\\
%     &= \int_{\mathbb{R}} g(x+{ka})g(x+{k'a}) e^{-2 \pi i ({j'b}-{jb}) x} dx \\
%     &= \mathcal{F}\{g(x+{ka})g(x+{k'a})\}({j'b}-{jb}),
% \end{align*}
% where we observed that the integral is just the Fourier transform of the product $g(x+{ka})g(x+{k'a})$ evaluated at $\xi = {j'b}-{jb}$. It is also the Fourier transform of a shifted product of B-splines centered at the origin $g(x+{ka})g(x-{ka})$, which can be seen by simply factoring the integrand. Without loss of generality, $|{ka}| < |{k'a}|$, we have a parabola centered at $x=\frac{{ka}+{k'a}}{2}$:
% \begin{align*}
%     g(x+{ka})g(x+{k'a}) &= (x-{ka}+1) (-x+{k'a}+1) \\
%     &= -x^2 +({ka}+{k'a})x + ({ka}{k'a} + {k'a} - {ka} +1).
% \end{align*}
% When ${ka}=-{k'a}$, the parabola is centered at $x=0$ and the integrand is ${g(x+{ka})g(x-{ka})}$. Using the time-frequency shift property of the Fourier transform, we then have
% \begin{align*}
%     \mathcal{F}\{g(x+{ka})g(x+{k'a})\}({j'b}-{jb})
%     = e^{2 \pi i (\frac{{ka} + {k'a}}{2})}  \mathcal{F}\{g(x+{ka})g(x-{ka})\}({j'b}-{jb})
% \end{align*}
% We also have that shifting the integrand in the other direction, $g(x-{ka})g(x-{k'a})$, conjugates the above:
% \begin{align*}
%     \mathcal{F}\{g(x-{ka})g(x-{k'a})\}({j'b}-{jb})
%     &= e^{-2 \pi i (\frac{{ka} + {k'a}}{2})}  \mathcal{F}\{g(x+{ka})g(x-{k'a})\}({j'b}-{jb}) 
% \end{align*}
% Therefore,
% \begin{align*}
%     (G_n^{[\ell]})_{n-k+1, n-k'+1} &= \mathcal{F}\{g(x+{ka})g(x+{k'a})\}({j'b}-{jb}) \\
%     &= \overline{\mathcal{F}\{g(x-{ka})g(x-{k'a})\}}({j'b}-{jb}) \\
%     &= \overline{(G_n^{[\ell]})_{k,k'}}
% \end{align*}


% \end{proof}





    


\section{Spectral analysis of sub-blocks of $G_n$ and $\mathscr{G}$}\label{sec:asymptotic}
%There are only a few cases of $a$ for which $\gamma_1\to 0$ and you can prove something about the infinite case based on interlacing.



We note that there is a Szeg\"{o} theorem for block Toeplitz matrices \cite{miranda2000asymptotic, gutierrez2008asymptotically, gutierrez2012block} that gives asymptotic results for the infinite block-Toeplitz matrix with blocks of size $k\times k$. Here, the sequence of block-Toeplitz matrices $G_n$ increases in both the number of blocks and the size $k$ of each block.

In this section we use the theory of Toeplitz matrices to describe the spectral properties of the block families
\[
T_n^{[\ell]}\in\mathbb{C}^{n\times n} \quad\text{and}\quad 
T_{\infty}^{[\ell]}\in\mathbb{C}^{\mathbb{Z}\times\mathbb{Z}},
\]
which govern the finite and infinite matrices $G_n^{[\ell]}$ and $\mathscr{G}^{[\ell]}$ through the  similarity relations proved in Section \ref{sec:gram_block}:
\[
G_n^{[\ell]} = D_n^{[\ell]} T_n^{[\ell]} D_n^{[\ell]}, 
\qquad 
\mathscr{G}^{[\ell]} = D_{\infty}^{[\ell]} T_{\infty}^{[\ell]} D_{\infty}^{[\ell]},
\]
where $D_n^{[\ell]}$ and $D_{\infty}^{[\ell]}$ are unitary diagonal matrices with unimodular entries, $(T_n^{[\ell]})_{k,k'}= t^{[\ell]}_{\,k-k'}$ for $k,k'\in \mathcal{I}_n$, and $(T_{\infty}^{[\ell]})_{k,k'}=t^{[\ell]}_{\,k-k'}$ for $k,k'\in \mathbb{Z}$.
Consequently, $G_n^{[\ell]}$ (resp. $\mathscr{G}^{[\ell]}$) and 
$T_n^{[\ell]}$ (resp. $T_{\infty}^{[\ell]}$) are unitarily equivalent and
share the same singular values (and the same spectra when the matrices are Hermitian).
Therefore, the asymptotic spectral behavior of $G_n^{[\ell]}$ (resp. $\mathscr{G}^{[\ell]}$)
is determined by the Toeplitz blocks $T_n^{[\ell]}$ (resp. $T_{\infty}^{[\ell]}$).










\subsection{Classical results}
We summarize classical results on Toeplitz matrices and refer to \cite{gray2006toeplitz, MR36936, zhu2017asymptotic, bottcher2000toeplitz} for comprehensive treatments.

Let \( \mathbf{S}^1 \coloneqq \{ z \in \mathbb{C} : |z| = 1 \} \) denote the complex unit circle. The \emph{Wiener algebra} \( W \) consists of functions \( f : \mathbf{S}^1 \to \mathbb{C} \) whose Fourier series converge absolutely:
\[
f(z) = \sum_{n=-\infty}^{\infty} \hat{f}_n z^n, \quad \|f\|_W = \sum_{n=-\infty}^{\infty} |\hat{f}_n| < \infty.
\]
The Fourier coefficients \( \hat{f}_n \) are given by
\[
\hat{f}_n = \frac{1}{2\pi} \int_0^{2\pi} f(e^{i\theta}) e^{-in\theta} \, d\theta.
\]
To each \( f \in W \) corresponds an infinite Toeplitz matrix \( T(f) \) defined by
\[
T(f) = \begin{bmatrix}
\hat{f}_0 & \hat{f}_{-1} & \hat{f}_{-2} & \cdots \\
\hat{f}_1 & \hat{f}_0 & \hat{f}_{-1} & \cdots \\
\hat{f}_2 & \hat{f}_1 & \hat{f}_0 & \cdots \\
\vdots & \vdots & \vdots & \ddots
\end{bmatrix}.
\]
The function \( f \), called the \emph{symbol} of the Toeplitz operator \( T(f) \), defines a bounded linear operator on \( \ell^p(\mathbb{Z}_+) \) for any \( 1 \leq p \leq \infty \), with operator norm satisfying \( \|T(f)\|_{\ell^p \to \ell^p} \leq \|f\|_W \) \cite{bottcher2005spectral}.

% The spectrum of \( T(f) \) is given by
% \[
% \operatorname{sp}(T(f)) 
% = f(\mathbf{S}^1) 
% \cup \left\{ \lambda \in \mathbb{C} \setminus f(\mathbf{S}^1) : \operatorname{wind}(f - \lambda) \neq 0 \right\},
% \]
% where \( \operatorname{wind}(f - \lambda) \) denotes the winding number of \( f - \lambda \) around the origin.

When the Fourier series of \( f \) has only finitely many non-zero coefficients, the associated Toeplitz matrix \( T(f) \) is a symmetric, banded bi-infinite matrix. In this case, \( f \) is called a \emph{Laurent polynomial}, and has the form
\[
f(z) = \sum_{j=-r}^{s} \hat{f}_j z^j, \quad z \in \mathbf{S}^1,
\]
for some integers \( r, s \geq 0 \) and real coefficients \( \hat{f}_j \in \mathbb{R} \).

In this case, \( T(f) \) has no eigenvalues, and its spectrum is the interval \( [\min f, \max f] \) \cite[Theorem~1.4]{bottcher2000toeplitz}.

\begin{theorem}[\cite{MR36936}]\label{thm:Tk_spec}
Let \( f \) be a real Laurent polynomial and let \( T(f) \) be the associated Toeplitz operator on \( \ell^2(\mathbb{Z}) \). Then
\[
\operatorname{sp}(T(f)) = \left[ \min_{z \in \mathbf{S}^1} f(z), \, \max_{z \in \mathbf{S}^1} f(z) \right].
\]
Moreover, if \( f \) is nonconstant, then \( T(f) \) has empty point spectrum (i.e., no eigenvalues).
\end{theorem}


\begin{theorem}[Lemma 4.1 in \cite{gray2006toeplitz} and Corollary 2.18 in \cite{bottcher2000toeplitz}]\label{cor-bdsfiniteteop}
Let \( T_n \) be the \( n \times n \) Hermitian Toeplitz matrix associated with the real-valued, continuous symbol \( t(x) \) defined on \( \mathbf{S}^1 \). Then the eigenvalues of \( T_n \) lie in the interval
\[
[\inf_{x} t(x), \sup_{x} t(x)],
\]
and satisfy the limits
\[
\lim_{n \to \infty} \lambda_1(T_n) = \inf_{x} t(x), \quad
\lim_{n \to \infty} \lambda_n(T_n) = \sup_{x} t(x).
\]
\end{theorem}

The following Szeg\"{o}-type limit theorem can be applied.
\begin{theorem}[Theorem 4.1, \cite{gray2006toeplitz}]\label{thm:Tn_lim_eigs}
Let \( T_n \) be a banded Toeplitz matrix with eigenvalues $\lambda_m(T_n)$ and real valued symbol $t$. Then for any function $F(x)$ continuous on $[\inf_xt(x), \sup_xt(x)]$, 

\[\lim_{n\to\infty} \frac{1}{n}\sum_{m=0}^{n-1}F(\lambda_m(T_n))=\frac{1}{2\pi}\int_{0}^{2\pi}F(t(\omega))d\omega. \]


\end{theorem}


\subsection{Spectral analysis of the infinite Toeplitz block $T_{\infty}^{[\ell]}$}

We fix a modulation difference $\ell$ in this section.



Let $T_{\infty}^{[\ell]}$ be the bi-infinite Toeplitz operator on
$\ell^2(\mathbb{Z})$ with entries $(T_{\infty}^{[\ell]})_{k,k'}=t^{[\ell]}_{\,k-k'}$,
where the coefficients $t^{[\ell]}_{\,j}$ are given by \eqref{eq:tjl}
and vanish for $|j|>\lfloor N/a\rfloor$. According to Corollary~\ref{cor:infBlocksTn}, $T_{\infty}^{[\ell]}$, is symmetric Toeplitz matrix with real entries. Its Laurent symbol is the real
trigonometric polynomial
\[
t^{[\ell]}(x)=\sum_{k=-\lfloor N/a\rfloor}^{\lfloor N/a\rfloor} t^{[\ell]}_{\,k}\,\cos(2\pi k x),
\qquad x\in[-\tfrac{1}{2},\tfrac{1}{2}].
\]
Figure \ref{fig:laurent_poly} shows plots of these Laurent polynomials for varying $a,b$ and $\ell$. In the case where the Toeplitz matrices are generated by an $N$th order B-spline, these Laurent polynomials have a closed analytic form.





\begin{figure} 
    \centering
   
         \includegraphics[width=\textwidth]{Figures/LaPoly.png}
         \caption{The Laurent polynomials $t^{[\ell]}(x) \in W$ associated to the Toeplitz blocks $T_{\infty}^{[\ell]} \in \mathcal{B}(l^p(\mathbb{Z}))$, $\ell=0, \hdots, 4$. The spectrum of $T_{\infty}^{[\ell]}$ is the interval $[\min t^{[\ell]}(x), \max t^{[\ell]}(x)]$. Here, $a=.3, .4, .5$ and the shades of blue correspond to different values of $b\in [1.5, 1.98]$.}
        \label{fig:laurent_poly}
  
\end{figure}






\begin{theorem} \label{thm:tl_sinc}
Let $N\geq2$. The Laurent polynomials $t^{[\ell]}(x)$ generated by the $N$th order
  B-spline $g=s_N$ have the form
  \begin{equation}
  t^{[\ell]}(x) = \frac{1}{a} \sum_{r \in \mathbb{Z}}
    \sinc^N\!\left(\frac{r - (x - ab\ell/2)}{a}\right)
	\sinc^N\!\left(\frac{r - (x + ab\ell/2)}{a}\right)
	\label{eq:tl_sinc}
  \end{equation}
\end{theorem}

  \begin{proof}  
% \begin{align*}
% t^{[\ell]}_k 
% &= \int_{\mathbb{R}} 
%     g\!\left(x - \frac{a k}{2}\right)
%     g\!\left(x + \frac{a k}{2}\right)
%     e^{2\pi i b \ell x}\,dx, \\[4pt]
% t^{[\ell]}(x)
% &= \sum_{k\in\mathbb{Z}} t^{[\ell]}_k\, e^{2\pi i k x}.
% \end{align*}

Applying the change of variables \(y = x - \tfrac{a k}{2}\), so \(x = y + \tfrac{a k}{2}\) to \eqref{eq:tjl}, we obtain
\begin{equation*}
t^{[\ell]}_k 
= e^{\pi i a b \ell k}
  \int_{\mathbb{R}} g(y)\, g(y + a k)\, e^{2\pi i b \ell y}\,dy.
\end{equation*}
Define \(f_\ell(y) = g(y)\,e^{2\pi i b \ell y}\).
Then
\begin{equation*}
t^{[\ell]}_k 
= e^{\pi i a b \ell k}\, (f_\ell * g)(a k),
\qquad 
(f_\ell * g)(z) = \int g(y)\, g(y+z)\, e^{2\pi i b \ell y}\,dy.
\end{equation*}

Hence the Laurent polynomial has the form
\begin{equation*}
t^{[\ell]}(x)
= \sum_{|k|\leq \lfloor \frac{N}{a} \rfloor} e^{\pi i a b \ell k}\, (f_\ell * g)(a k)\, e^{2\pi i k x},
\end{equation*}
% The Poisson summation formula states for any Schwartz function \(f\),
% \[
% \sum_{k\in\mathbb{Z}} f(a k)\, e^{2\pi i k x}
% = \frac{1}{a} \sum_{r\in\mathbb{Z}} 
%   \widehat{f}\!\left(\frac{r-x}{a}\right),
% \quad
% \widehat{f}(\xi) = \int f(y)\,e^{-2\pi i y\xi}\,dy.
% \]
and applying the Poisson summation formula gives
\begin{equation*}
t^{[\ell]}(x)
= \frac{1}{a}\sum_{r\in\mathbb{Z}}
  \widehat{(f_\ell * g)}\!\left(\frac{r - (x - a b \ell/2)}{a}\right)= \frac{1}{a}\sum_{r\in\mathbb{Z}}
   \widehat{g}\!\left(\tfrac{r - (x - a b \ell/2)}{a}\right)
   \widehat{g}\!\left(\tfrac{r - (x - a b \ell/2)}{a} - b\ell\right).
\end{equation*}
% The Fourier transform of each term in the summand is given by \[
% \widehat{(f_\ell * g)}(\xi)
% = \widehat{f_\ell}(\xi)\,\widehat{g}(\xi),
% \text{ where }
% \widehat{f_\ell}(\xi)
% = \widehat{g}(\xi - b\ell).
% \]
% Therefore,
% \[
% t^{[\ell]}(x)
% = \frac{1}{a}\sum_{r\in\mathbb{Z}}
%    \widehat{g}\!\left(\tfrac{r - (x - a b \ell/2)}{a}\right)
%    \widehat{g}\!\left(\tfrac{r - (x - a b \ell/2)}{a} - b\ell\right).
% \]
The Fourier transform of $g=s_N$ has the form
\begin{equation*}
\widehat{g}(\xi)
= \sinc^N(\xi)
= \left(\frac{\sin(\pi \xi)}{\pi \xi}\right)^N.
\end{equation*}
Hence the representation of $t^{[\ell]}(x)$ is given by
\eqref{eq:tl_sinc}. 

\end{proof}

%\begin{corollary}
%The Laurent polynomials \(t^{[\ell]}(x)\) generated by the $n$th order
%B-spline $g = s_n$ have the form
%\[
%t^{[\ell]}(x) = \frac{1}{a} \sum_{r \in \mathbb{Z}}
%\sinc^n \left(\frac{r - (x - ab\ell/2)}{a}\right)
%\sinc^n \left(\frac{r - (x - ab\ell/2)}{a} - b\ell\right)
%\]
%\end{corollary}
%
%\begin{proof}
%The proof is the same as above, replacing $s_2$ with $s_n$ and noting
%that
%\[
%s_n = s_{n-1} * s_1
%\]
%from which it follows that
%\[
%\widehat{g}(\xi) = \sinc^n(\xi)
%\]
%\end{proof}


\begin{corollary} \label{corr: laurent_max_min}
    Suppose that $a = \tfrac{1}{p}$ with $p\in\mathbb{N}$,  $p\ge 2$, and $\ell\ge 0$ and let
$T_{\infty}^{[\ell]}$ denote the Toeplitz operator on $\ell^2(\mathbb{Z})$ with symbol
$t^{[\ell]}(x)$. Then
\begin{equation*}
\Spec\big(T_{\infty}^{[\ell]}\big) \subset [0,p].
\end{equation*}
Moreover, when $\ell=0$ we have 

\begin{equation*}
\Spec\big(T_{\infty}^{[0]}\big)= [0,p].
\end{equation*}

    \end{corollary}

    \begin{proof}
By Theorem~\ref{thm:tl_sinc} with $\ell=0$ and $a = 1/p$, we have
\begin{equation*}
t^{[0]}(x)
= \frac{1}{a}\sum_{r\in\mathbb{Z}} \sinc^{2N}\!\left(\frac{r-x}{a}\right)
= p \sum_{r\in\mathbb{Z}} \sinc^{2N}\big(p(r-x)\big).
\end{equation*}
At integer points $x\in\mathbb{Z}$, only the term $r=x$ contributes, so $t^{[0]}(x)=p$.

On the other hand, if $x = k/p$ with $k\in\mathbb{Z}\setminus p\mathbb{Z}$, then
$p(r-x) = pr-k$ is a nonzero integer for every $r\in\mathbb{Z}$, and hence
each $\sinc$ factor vanishes. Thus $t^{[0]}(k/p)=0$ for
$k\in\mathbb{Z}\setminus p\mathbb{Z}$.
Since $t^{[0]}$ is continuous and nonnegative, it follows that by Theorem \ref{thm:Tk_spec}
\begin{equation*}
\Spec(T_{\infty}^{[0]})=\left[\inf_x t^{[0]}(x), \sup_x t^{[0]}(x)\right] = [0,p].
\end{equation*}






\end{proof}



The numerical results in Figure \ref{fig:Tijdecay} demonstrate $N$th order decay of the spectral width $|\sup(t^{[\ell]}(x))-\inf(t^{[\ell]}(x))|$ of the finite Gram matrices $G_n$ for Gabor systems generated by the $N$th order B-spline. This gives a notion of block diagonal dominance of the matrix $G_n$.
\begin{figure}[h!]
	\centering
    \includegraphics[width=\linewidth]{Figures/tndecay.23_1.7.png}
	\caption{Width of spectrum of $t^{[\ell]}(x)$ for $a=.23$, $b=1.7$ corresponding to the $N$th order B-spline. The dashed lines show the corresponding $N$th order decay proved in Lemma \ref{lem:decay}.}\label{fig:Tijdecay}
\end{figure}

\begin{lemma}\label{lem:decay}
    Let $|\ell|>0$ and let $t^{[\ell]}(x)$ be the Laurent polynomial corresponding to the $N$th-order B-spline. Then, there exists a $C>0$ such that
    \[|\sup(t^{[\ell]}(x))-\inf(t^{[\ell]}(x))|\leq C {\ell}^{-N}.\]
\end{lemma}

\begin{proof} Suppose that $|\ell|>\tfrac{4}{ab}$ and let $x\in [-1/2, 1/2]$.  From~\eqref{eq:tl_sinc} we have:

\begin{align*}
|t^{[\ell]}(x)|& \leq \tfrac{a^{2N-1}}{\pi^{2N}} \tfrac{1}{|x-ab\ell/2|^N |x+ ab\ell/2|^N} + \tfrac{a^{2N-1}}{\pi^{2N}} \sum_{r=1}^{\infty} \tfrac{1}{|r-(x-ab\ell/2)|^N |r-(x+ab\ell/2)|^N} + \tfrac{a^{2N-1}}{\pi^{2N}} \sum_{r=-\infty}^{-1} \tfrac{1}{|r-(x-ab\ell/2)|^N |r-(x+ab\ell/2)|^N}\\
& = \tfrac{a^{2N-1}}{\pi^{2N}} \tfrac{1}{|x-ab\ell/2|^N |x+ ab\ell/2|^N} + \tfrac{a^{2N-1}}{\pi^{2N}} \sum_{r=1}^{\infty} \tfrac{1}{|r-(x-ab\ell/2)|^N |r-(x+ab\ell/2)|^N} + \tfrac{a^{2N-1}}{\pi^{2N}} \sum_{r=1}^{\infty} \tfrac{1}{|r+(x-ab\ell/2)|^N |r+(x+ab\ell/2)|^N}.
\end{align*}

We assume that $\ell>\tfrac{4}{ab}$ and the case $\ell<-\tfrac{4}{ab}$ is treated similarly.  

It follows that $$\sup_{x\in [-1/2, 1/2]}\tfrac{1}{|x-ab\ell/2|^N |x+ ab\ell/2|^N} \leq 4^{3N}(ab)^{-2N}\  \ell^{-2N}. $$ 

Let $L_0=\lfloor \tfrac{ab\ell +1}{2} \rfloor$. Then,


\begin{align*}
\sum_{r=1}^{\infty} \tfrac{1}{|r-(x-ab\ell/2)|^N |r-(x+ab\ell/2)|^N} &= \sum_{r=1}^{L_0-1} \tfrac{1}{|r-(x-ab\ell/2)|^N |r-(x+ab\ell/2)|^N} + \sum_{r=L_0}^{\infty} \tfrac{1}{|r-(x-ab\ell/2)|^N |r-(x+ab\ell/2)|^N}.
\end{align*}

By writing $\tfrac{ab\ell +1}{2}=L_0+\theta$
for some $\theta \in (0, 1)$ we see that 
\begin{align*}
\sum_{r=L_0}^{\infty} \tfrac{1}{|r-(x-ab\ell/2)|^N |r-(x+ab\ell/2)|^N} &\leq \sum_{r=L_0}^\infty \tfrac{1}{|r+\tfrac{ab\ell-1}{2}|^N} \tfrac{1}{|r-\tfrac{ab\ell +1}{2}|^N}\\
&\leq (L_0 + \tfrac{ab\ell-1}{2})^{-N} \sum_{r=L_0}^\infty \tfrac{1}{|r-\tfrac{ab\ell +1}{2}|^N}\\
&\leq 2^N(ab)^{-N}\ell^{-N} \sum_{r=L_0}^\infty \tfrac{1}{|r-L_0-\theta |^N}\\
&= 2^N(ab)^{-N}\ell^{-N} \sum_{r=0}^\infty \tfrac{1}{|r-\theta |^N}\leq_N \ell^{-N}.
\end{align*}

Next, for each $r\in \{1, 2, \hdots, L_0-1\}$ we have that $$\sup_{x\in [-1/2, 1/2]} \tfrac{1}{|r-(x-ab\ell/2)|^N} \leq \tfrac{2^N}{(1+ab\ell)^N}\leq 2^N(ab)^{-N}\ell^{-N}.$$

Moreover, $$\sup_{x\in [-1/2, 1/2]} \sum_{r=1}^{L_0-1} \tfrac{1}{ |r-(x+ab\ell/2)|^N}= \sup_{x\in [-1/2, 1/2]} \sum_{r=1}^{L_0-1} \tfrac{1}{ |r-x -L_0 -\theta +1/2|^N}= \sup_{x\in [0, 1]} \sum_{r=1}^{L_0-1} \tfrac{1}{ |r -L_0 +x -\theta |^N}<\infty $$ because $r-L_0-\theta \leq -1-\theta<-1$. 

We can estimate $\sum_{r=1}^{\infty} \tfrac{1}{|r+(x-ab\ell/2)|^N |r+(x+ab\ell/2)|^N} $ similarly to conclude that 




$$|t^{[\ell]}(x)| \leq a^{2N-1}\pi^{-2N} (4^{3N}(ab)^{-2N}\ell^{-2N} +C_12^{N+1}(ab)^{-N}\ell^{-N}+C_22^{N+1}(ab)^{-N}\ell^{-N})\leq C_N \ell^{-N}$$ for some constants $C_1, C_2$, and $C_N$. 
    
\end{proof}





















\subsection{Spectral analysis of the finite Gram matrix $G_n$}


Given Corollary~\ref{cor:infBlocksTn}, the infinite block $\mathscr{G}^{[\ell]}$  within $\mathscr{G}$  has the same singular values as the real-valued, banded Toeplitz matrix $T_{\infty}^{[\ell]}$ whose entries are given by \eqref{eq:tjl}. 


Since $G_n^{[0]}$ is a principal sub-matrix of $G_n$, Cauchy's interlacing theorem states that their eigenvalues interlace, leading to bounds  on the minimum and maximum eigenvalues $\lambda_1(G_n)$ and  $\lambda_{n^2}(G_n)$ of ${G}_n$: 
%\begin{remark} Using the classical  Szego-type limit theorems \cite{gray2006toeplitz} we compare the extreme values of the spectrum of $T_n^{[\ell]}$ to that of $T_{\infty}^{[\ell]}$.
%
%\begin{enumerate}
   % \item The minimum and maximum eigenvalues of the $n\times n$ Toeplitz matrix $T_n^{[\ell]}$ belong to the interval $[\inf_x t^{[\ell]}(x), \sup_x t^{[\ell]}(x)]$ with
   % \begin{align*}
     %   & \lim_{n\rightarrow \infty} \lambda_1(T_n^{[\ell]}) = \inf_x t^{[\ell]}(x) \\
     %   & \lim_{n\rightarrow \infty} \lambda_n(T_n^{[\ell]}) = \sup_x t^{[\ell]}(x)
   % \end{align*}
    %\item 
    %Since blocks $G_n^{[0]}$ are principal sub-matrices and hence the eigenvalues interlace with ${G}_n$,
    \begin{align*}
        \lambda_k(G_n) \leq \lambda_k(G_n^{[0]}) \leq \lambda_{k+n^2-n}(G_n), \qquad k=1,\hdots n.
    \end{align*}

    
           In particular,  $\lambda_1(G_n) \leq \lambda_1(G_n^{[0]})$ and $\lambda_n(G_n^{[0]}) \leq  \lambda_{n^2}(G_n)$. Figure \ref{fig:minmax_Gn_eigns} demonstrates this for $n=15$.
%\end{enumerate}   
%\end{remark}




Consequently, we have the following result.

%Note that these blocks on the diagonal $\Gamma^{[0]}$ correspond to sampling the STFT along lines of the lattice $\Lambda$.

% Let $t^{[0]}\in W$ be the Laurent polynomial $$t^{[0]}(x) = 2/3 + \sum_{k=1}^{\lfloor 1/a \rfloor} t_j \cos(2 \pi j x) + \frac{1}{3}\sum_{j=\lfloor 1/a \rfloor + 1}^{\lfloor 2/a \rfloor} t_j \cos(2 \pi j x),$$  with the coefficients given as  the inner products $\langle g_{j,k}, g_{j',k'}\rangle_{L^2(\mathbb{R})}$ on the Toeplitz bands of $\Gamma^{[0]}$ indexed by $j=|k-k'|$. 

% These coefficients are explicitly given by 
%     \begin{equation*}
%         t_j = \begin{cases}
%             (aj)^3-2(aj)^2+\frac{4}{3}, \ j \leq \lfloor \frac{1}{a}\rfloor  \\
%             (2-aj)^3, \ \lfloor \frac{1}{a} \rfloor \leq j \leq \lfloor \frac{2}{a} \rfloor \\
%             0, \ j \geq \lfloor \frac{2}{a} \rfloor
%         \end{cases}
%     \end{equation*}
    

    
% {\color{blue} combine this with the next one or delete
% \begin{theorem} \label{thm: gii_spec} The infinite block $\Gamma^{[0]} \in \mathcal{B}(\ell^2(\mathbb{Z}_+))$  on the diagonal of $\Gamma$ is precisely the Toeplitz matrix $\cT(t^{[0]})$. In addition, $\Gamma^{[0]}$ has no eigenvalues and its spectrum is the interval
%     \begin{equation*}
%         \text{Spec} \big(\Gamma^{[0]}\big) = [\inf_x t^{[0]}(x), \sup_x t^{[0]}(x)].
%     \end{equation*}
% \end{theorem}

% }

	
%    By Theorem \ref{thm: diag_blocks}, the $\Gamma^{[0]}$ are Toeplitz, symmetric, real-valued, and banded. They therefore have no eigenvalues and the spectrum is given by a Laurent polynomial constructed from the bands of the $\Gamma^{[0]}$. The entries $G_n^{[0]}_{k,k'}$ for $j,j' \in \mathbb{Z_+}$ are given by integrating one shifted B-spline against another as seen in Equation \ref{eq: gram_entry} and has piecewise form. If ${ka} \leq {k'a}$, then:
%    \begin{equation*}
 %       (G_n^{[\ell]})_{k,k'} = 
  %      \begin{cases}
            %\int_{{k'a}-1}^{{ka}}(x+1-{ka})(x+1-{k'a}) dx \\
          %  + \int_{{ka}}^{{k'a}} (x-{k'a}+1)(-x+{ka}+1) dx \\
         %   + \int_{{k'a}}^{{ka}+1}(-x+1+{ka})(-x+1+{k'a})dx 
         %   & 0 \leq {ka}-{k'a} < 1 \\
            %\int_{{ka}'-1}^{{ka}+1} (x-{k'a}+1)(-x+{ka}+1) & 1 \leq {ka}-{k'a} < 2 \\
           % 0 & 2 \leq {ka}-{k'a}
      %  \end{cases}
   % \end{equation*}
  %  If ${k'a} \leq {ka}$, then due to symmetry the entries $G_n^{[0]}_{k,k'}$ have the same form but with $j$ and $j'$ swapped. Since ${ka}=-a(\lfloor n/2\rfloor + (j-1))$, we have that ${ka}-{k'a} = a(j-j')$. As before, index the Toeplitz coefficients as $k = j-j'$. Then, solving these integrals yields:
   % \begin{equation*}
    %     (G_n^{[\ell]})_{k,k'} = G^{ii}_{j-j'}=G^{ii}_k =
    %     \begin{cases}
    %         (ak+2)^3/6 & -2/a \leq k \leq -1/a \\
    %         (ak+1)^2(-ak+2)/3
    %         - (ak)^3/6-ak(ak+1)
    %         &  -1/a \leq k \leq 0 \\
    %         (-ak+1)^2(ak+2)/3
    %         + (ak)^3/6+ak(-ak+1)
    %         &  0 \leq k \leq 1/a \\
    %         (-ak+2)^3/6 & 1/a \leq k \leq 2/a \\
    %         0 & |k| \geq 2/a
    %     \end{cases}
    % \end{equation*}

    % Because the coefficients for positive and negative index $l$ are equal, the Laurent polynomial $g^{ii}(x) \in W$ associated to $\Gamma^{[0]}$ is real-valued. The $g^{ii}(x)$ has the form:
    % \begin{align*}
    %     g^{ii}(x) &= \sum_{k=-\lfloor 2/a \rfloor}^{\lfloor 2/a \rfloor}G^{ii}_k e^{2 \pi i k x} \\
    %     &= 2/3 + 2\sum_{k=1}^{\lfloor 2/a \rfloor}G_k^{ii} \cos(2 \pi k x) \\
    %     &= 2/3 + \sum_{k=1}^{\lfloor 1/a \rfloor} \big( (ak)^3 - 2(ak)^2 + \frac{4 }{3} \big) \cos(2 \pi k x)  + \frac{1}{3}\sum_{k=\lfloor 1/a \rfloor + 1}^{\lfloor 2/a \rfloor} (2-ak)^3 \cos(2 \pi k x)
    % \end{align*}
    
    % Clearly, $\Gamma^{[0]} \in \mathcal{B}(\ell^p(\mathbb{Z}_+))$ as the Fourier series is finite with finite coefficients.


% For certain values of $a$ and $b$ more can be said about the spectrum of $T_{\infty}^{[0]}$. 

%\begin{example}
% For example, if we take $a=1/4$ and $b=2$ (known to \textit{not} be in the frame set for the 2nd-order B-spline), then $g^{ii}(x)$ (Figure \ref{fig: theo_gii}) has the form:
% \begin{align*}
%     g^{ii}(x) &= 2/3 + \sum_{k=1}^{4} (\frac{k^3}{64} - \frac{k^2}{8} + \frac{4}{3})\cos(2 \pi kx) +
%     \frac{1}{3} \sum_{k=5}^8 (2-\frac{k}{4})^3 \cos(2 \pi k x)
% \end{align*}
% As all the coefficients are positive, the $g^{ii}(x)$ attains its maximum when $\cos(2\pi kx)=1$. Using a standard counting argument we can explicitly compute the maximum and minimum of $g^{ii}(x)$ in certain cases. 

% Recall the formulas for the sums of natural numbers, their squares, and cubes
% \begin{align*}
%     \sum_{k=1}^n k &= \frac{n(n+1)}{2} \\
%     \sum_{k=1}^n k^2 &= \frac{n(n+1)(2n+1)}{6} \\
%     \sum_{k=1} k^3 &= \Big(\frac{n(n+1)}{2}\Big)^2
% \end{align*}
% and the alternating sums of natural numbers, their squares, and cubes
% \begin{align*}
%     \sum_{k=1}^n(-1)^k k &= \begin{cases}
%         -(\frac{n+1}{2})^2 + (\frac{n-1}{2})(\frac{n-1}{2}+1),  \ n \text{ even}\\
%         -(\frac{n}{2})^2 - (\frac{n}{2})(\frac{n}{2}+1), \ n \text{ odd}
%     \end{cases} \\
%     \sum_{k=1}^n (-1)^kk^2 &= \begin{cases}
%         \frac{n(n+1)}{2}, \ n \text{ even} \\
%         -\frac{n(n+1)}{2}, \ n \text{ odd}
%     \end{cases} \\
%     \sum_{k=1}^n (-1)^k k^3 &= \begin{cases}
%         \frac{2n^3 + 3n^2 -1}{4}, \ n \text{ even} \\
%         \frac{-n^2(2n+3)}{4}, \ n \text{ odd}.
%     \end{cases}
% \end{align*}
%\end{example}





%\begin{proof}
    % Without loss of generality, take $k=0$. To simplify notation, let $A \coloneqq \lfloor \frac{1}{a}\rfloor$ and $B \coloneqq \lfloor \frac{2}{a} \rfloor = 2A$. Note, the interesting cases are when $a <1$ and there greater than zero or one term in the summation. Then,
    
    % Apply the formulas for the sum of natural numbers, their squares, and cubes. The above becomes:
    % \begin{align*}
    %      t^{[0]}(0) &= 2/3 + a^3 \big (\frac{A^2(A+1)^2}{4} \big) - 2a^2 \big(\frac{A(A+1)(2A+1)}{6}  \big) + \frac{4A}{3} \\
    %     &+ \frac{1}{3} \big(-a^3 (\frac{B^2(B+1)^2}{4}) + 6a^2 (\frac{B(B+1)(2B+1)}{6}) -12a (\frac{B(B+1)}{2}) + 8B \\
    %     &+a^3(\frac{A^2(A+1)^2}{4}) - 6a^2(\frac{A(A+1)(2A+1)}{6}) +12a(\frac{A(A+1)}{2}) - 8A\big)
    % \end{align*}
    % Simplifying:
    % \begin{align*}
    %     t^{[0]}(0) &= 2/3 + \frac{a^3A^4}{4} + \frac{a^3A^3}{2} + \frac{a^3A^2}{4} - \frac{2a^2A^3}{3} - a^2A^2-\frac{a^2A}{3} + \frac{4A}{3} \\
    %     &- \frac{1}{12}a^3B^4 - \frac{a^3B^3}{6} - \frac{a^3B^2}{12}+\frac{2a^2B^3}{3} + a^2B^2 + \frac{a^2B}{3} - 2aB^2 -2aB + \frac{8B}{3} \\
    %     &+\frac{1}{12}a^3A^4 + \frac{a^3A^3}{6} + \frac{a^3A^2}{12} - \frac{2a^2 A^3}{3} - a^2A^2 - \frac{a^2A}{3} + 2aA^2 + 2aA - \frac{8A}{3}
    % \end{align*}
    % Recalling $A=\frac{1}{a}$ and $B=2A$ by assumption, combining like terms and simplifying, we have
    % \begin{equation*}
    %     t^{[0]}(0) = \frac{1}{a}.
    % \end{equation*}
%     To compute the minimum without loss of generality, take $x=1/2$. Then,
%     \begin{align*}
%         t^{[0]}(1/2) &= 2/3 + \sum_{k=1}^{A} (-1)^k (ak)^3 - 2\sum_{k=1}^{A}(-1)^k(ak)^2 + \sum_{k=1}^A (-1)^k \frac{4}{3} \\
%         &+ \frac{1}{3} \big(\sum_{k=1}^{B} (-1)^k (2-ak)^3 - \sum_{k= 1}^{A} (-1)^k(2-ak)^3 \big) \\
%         &= 2/3 +  a^3\sum_{k=1}^{A} (-1)^k k^3 - 2a^2\sum_{k=1}^{A} (-1)^k k^2 + \sum_{k=1}^A (-1)^k \frac{4}{3} \\
%         &+ \frac{1}{3} \big(-a^3\sum_{k=1}^{B}(-1)^k k^3 + 6a^2 \sum_{k=1}^{B} (-1)^k k^2 -12a \sum_{k=1}^{B} (-1)^k k + \sum_{k=1}^B (-1)^k8 \\
%         &- (-a^3\sum_{k=1}^A (-1)^k k^3 + 6a^2 \sum_{k=1}^A (-1)^k k^2 -12a \sum_{k=1}^A (-1)^k k + \sum_{k=1}^A (-1)^k 8) \big)
%     \end{align*}

%     Apply the formulas for alternating sums of natural numbers, their squares, and cubes. These formulas depend even and odd cases. When $A=\frac{1}{a}$ is odd, we have:
%     \begin{align*}
%         t^{[0]}(1/2) &= 2a(\frac{2}{a} + 1) + \frac{\big(a(\frac{2}{a} + 3) \big)}{3} - \frac{\big( a(\frac{4}{a}+3)\big)}{3} - 4a \big(\frac{\frac{1}{a}+1}{a}- \frac{1}{a^2}) \\
%         &+ 4a \big(\frac{(\frac{1}{2a} + 1)}{2a}-\frac{1}{4a^2} \big) - 2a(\frac{1}{a} +1 ) + \frac{2}{3} \\
%         &= 0
%     \end{align*}
%     The computation for the even case proceeds similarly.
% \end{proof}





%Toeplitz structure in the finite Gram matrices ${G}_n$, we can understand the asymptotic behavior of its eigenvalues and spectrum. This is due to a Fourier series one can construct from the Toeplitz bands and a class of \textit{asymptotically equivalent} circulant matrices whose eigenvalues we explicitly provide. 

%In this section we focus on the spectrum of the finite and infinite blocks $G_n^{[0]}$ and $\Gamma^{[0]}$ for $i \in \{1,2, \ldots \}$ to make conclusions about the spectra of the finite and infinite Gram matrices ${G}_n$ and $G(a,b)$, respectively.


%-=--=\subsection{Spectral analysis of finite section of the Gram matrix}

%along with the  classical  Szego-type limit theorems \cite{gray2006toeplitz} we have the following result.























% \begin{corollary} \color{red} combine with previous \label{corr: upper_frame_bounds}
%      When $\lfloor \frac{1}{a} \rfloor = \frac{1}{a}$, the upper frame bound of the Gabor system $\cG(g,a\Z\times b\Z)$ is $B(a,b)=\frac{1}{a}$.
% \end{corollary}
% That $B(a,b)=\frac{1}{a}$ follows from the fact that the upper frame bound is known to be the operator norm $B(a,b)=||G(a,b)||$. 




% Analysis of the Laurent polynomial $t^{[0]}(x)$ including its maximum and minimum provides information on the spectrum of the infinite Gram matrix $G(a,b)$ and its finite counterpart ${G}_n$, and hence information on their frame bounds. This is summarized in Corollaries \ref{corr: upper_frame_bounds} and \ref{corr: eigs} and is due to:
% {\color{red} when a=1/p, previous results, otherwise, when you can't compute explicitly you can estimate it with 1. below and with 2. you get the benefit of estimates of frame bounds interlacing}




\begin{corollary} \label{corr: eigs}
   For $n\geq 1$, there exists an $K \in \mathbb{N}$ such that $\forall n \geq K$, we have the upper and lower estimates:
    \begin{align}
        & \lambda_1(G_n) \leq \inf_x t^{[0]}(x) \leq  \lambda_1(G_n^{[0]}) \\
        & \lambda_n(G_n^{[0]}) \leq \sup_x t^{[0]}(x) \leq \lambda_n(G_n).
    \end{align}
\end{corollary}

\begin{proof}
We will only establish the first set of inequalities as the second one is proved in a similar manner. 

We note that $\inf_x t^{[0]}(x) \leq  \lambda_1(T_n^{[0]}) $ follows from Corollary~\ref{cor-bdsfiniteteop}. Next, by the interlacing relation, we know that  $\lambda_1(G_n)\leq \lambda_1(G_n^{[0]})$. It follows that $$\limsup_{n}\lambda_1(G_n)\leq \limsup_{n} \lambda_1(G_n^{[0]})= \lim_{n} \lambda_1(G_n^{[0]})=\inf_x t^{[0]}(x). $$ Thus, there exists an integer $K\geq 1$ such that for all $n\geq K$, we have 
$$\lambda_1(G_n)\leq \inf_x t^{[0]}(x).$$



\end{proof}



\begin{figure}
    \centering
    \includegraphics[width=\linewidth]{Figures/min_eig.png}
    \includegraphics[width=\linewidth]{Figures/max_eig.png}
    \caption{Plots of $\lambda_1(G_n)$ (top left), $\min t^{[0]}(x)$ (top right), $\lambda_n(G_n)$ (bottom left), and $\max t^{[0]}(x)$ (bottom right) for $n=15$ demonstrate the results of Corollary \ref{corr: eigs}.}
    \label{fig:minmax_Gn_eigns}
    
\end{figure}

\subsection{Asymptotically Equivalent Circulant Matrices}\label{sec:circulant}

The exact computation of the spectrum of a Toeplitz matrix \( G_n \), especially for large \( n \), is generally intractable. However, Theorem~\ref{thm:offdiagTH} shows that the spectral behavior of \( G_n \) is governed by the singular values of the Toeplitz blocks \( T_n^{[\ell]} \). To analyze these efficiently, we introduce a family of \emph{circulant matrices} \( C_n^{[\ell]} \) that approximate the Toeplitz matrices \( T_n^{[\ell]} \) in a precise asymptotic sense.

\subsubsection{Background on Asymptotic Equivalence}

% Two sequences of matrices \( \{A_n\} \) and \( \{C_n\} \) are said to be \emph{asymptotically equivalent}, denoted \( A_n \sim C_n \), if:
% \begin{enumerate}
%     \item Both sequences are uniformly bounded in operator norm: \( \|A_n\|_2, \|C_n\|_2 \leq M \) for some constant \( M \) independent of \( n \),
%     \item The difference vanishes in the normalized Hilbert–Schmidt norm:
%     \[
%     \lim_{n \to \infty} \left( \frac{1}{n} \sum_{i,j=1}^n |(A_n)_{i,j} - (C_n)_{i,j}|^2 \right)^{1/2} = 0.
%     \]
% \end{enumerate}


Let \( T_n \) be a banded Toeplitz matrix of bandwidth \( m \), defined by a symbol \( t(x) = \sum_{k=-m}^{m} t_k e^{2\pi i k x} \). Define the corresponding circulant matrix \( C_n \in \mathbb{C}^{n \times n} \) by wrapping the Toeplitz coefficients as:
\[
c_k = 
\begin{cases}
t_{-k}, & 0 \leq k \leq m, \\
t_{n-k}, & n-m \leq k \leq n-1, \\
0, & \text{otherwise}.
\end{cases}
\]
This defines the first row of \( C_n \), from which the rest is filled by circular shifts.

\begin{lemma}[Lemma 4.2 in \cite{gray2006toeplitz}]\label{lem:circeigconv} The matrices $T_n$ and $C_n$ are asymptotically equivalent, denoted by $T_n\sim C_n$ i.e.,
\begin{enumerate}
    \item Both sequences are uniformly bounded in operator norm: \( \|T_n\|_2, \|C_n\|_2 \leq M \) for some constant \( M \) independent of \( n \),
    \item The difference vanishes in the normalized Hilbert–Schmidt norm:
    \[
    \lim_{n \to \infty} \left( \frac{1}{n} \sum_{i,j=1}^n |(T_n)_{i,j} - (C_n)_{i,j}|^2 \right)^{1/2} = 0.
    \]
\end{enumerate}
    
\end{lemma}

When $T_n\sim C_n$ holds for Hermitian matrices, their eigenvalues are said to be \emph{asymptotically equally distributed}, meaning that for any integer \( s \geq 1 \),
\[
\lim_{n \to \infty} \frac{1}{n} \sum_{m=1}^{n} \left( \lambda^s_m(A_n) - \lambda^s_m(C_n) \right) = 0.
\]
For general theory, see \cite[Chapter 2]{gray2006toeplitz}. In special cases, such as banded Toeplitz matrices, stronger pointwise convergence of eigenvalues is known \cite{zhu2017asymptotic}.


Circulant matrices admit explicit spectral decompositions. Specifically:

\begin{theorem}[Theorem 5.15, \cite{bottcher2000toeplitz}]\label{thm:C_eigs}
Let \( C_n \) be the circulant matrix constructed from a trigonometric polynomial \( t \). Then its eigenvalues are given (up to permutation) by uniformly spaced samples of \( t \):
\[
\operatorname{spec}(C_n) = \left\{ t\left( \tfrac{p}{n} \right) : p = 0, \dots, n-1 \right\}.
\]
\end{theorem}






\subsubsection{Application to Gram Matrices for $\mathcal{G}(s_N,\Lambda_n)$}


Let \( T_n^{[\ell]} \) denote the Toeplitz blocks with symbol $t^{\ell}(x)$ generated by the second order B-spline, and define \( C_n^{[\ell]} \) to be the circulant matrix associated with the same symbol. 

Therefore, letting \( \{ \psi_1, \dots, \psi_n \} \) be the ordered eigenvalues of \( C_n^{[\ell]} \), i.e.,
\[
\psi_m = t^{[\ell]}\left( \tfrac{p_m}{n} \right)
\quad \text{for } p_m \in \{0, \dots, n-1\},
\]
we have that for any \( s \geq 1 \),
\[
\lim_{n \to \infty} \frac{1}{n} \sum_{m=1}^{n} \left( \lambda^s_m(T_n^{[\ell]}) - \psi^s_m \right) = 0.
\]
This approximation enables efficient estimation of spectral bounds for the finite Gram matrix \( G_n \). The asymptotic behavior of the difference in eigenvalues between $C_n^{[\ell]}$ and $T_n^{[\ell]}$ is shown in Figure \ref{fig: eig_diff}.

\begin{figure}[h!]
    \centering
    \includegraphics[width=0.9\textwidth]{Figures/eigen_convergence.png}
    \caption{Lemma \ref{lem:circeigconv} gives the convergence of eigenvalues of the Toeplitz block \( T_n^{[\ell]} \) to those of the circulant matrix \( C_n^{[\ell]} \). Here, $T_n^{[\ell]}$ correspond to the subblocks $G_n^{[\ell]}$ of the Gram matrices $G_n$ for $\mathcal{G}(s_2,\Lambda_n)$.}
    \label{fig: eig_diff}
\end{figure}

%{\color{red} Conjecture: It appears that the max difference for $\ell=0$ may be independent of $N$, the order of the $B-$spline.}



% \subsection{Asymptotically equivalent circulant matrices}
% It can be rather difficult to compute exactly the spectrum of $G_n$, especially for large values of $n$. {However, by Theorem~\ref{thm:offdiagTH} the spectrum of $G_n^{[\ell]}$ consists of the singular values of the Toeplitz matrices $T_n^{[\ell]}$}. At the same time, a sequence of circulant matrices $C_n^{[\ell]}$ that are \emph{asymptotically equivalent} matrices to the Toeplitz matrices $T_n^{[\ell]}$ can be constructed. Because, the spectra of circulant matrices can be computed explicity, spectral information of $G_n^{[\ell]}$ can be approximated. In particular,  we will obtain estimates on the eigenvalues $\lambda_1(G_n)$ and $ \lambda_{n^2}(G_n) $. 

% \subsubsection{Background}

% {\color{red} Gray book chapter 4. section 4.}

% Recall that two sequences of matrices $\{A_n\}_{n=1}^{\infty}, \{C_n\}_{n=1}^{\infty}$ are said to be \textit{asymptotically equivalent}, $\{A_n\} \sim \{C_n\}$, if:
% \begin{enumerate}
%     \item $\{A_n\}_{n=1}^{\infty}$ and $\{C_n\}_{n=1}^{\infty}$ are uniformly bounded in the $||\cdot||_2 \coloneqq \sqrt{\lambda_{\text{max}}(A^*A)}$ norm, \begin{equation*}
%       \exists M<\infty \text{ such that }\forall n, ||A_n||_2\leq M \text{ and }  ||C_n||_2 \leq M 
%     \end{equation*}
%     \item $\lim_{n\to \infty}|A_n - C_n| = 0$, where $|\cdot|$ is the scaled Hilbert-Schmidt norm $|A| \coloneqq (\frac{1}{n} \sum_{i,j=1}^n |A_{i,j}|^2)^{1/2}$
% \end{enumerate}
% Asymptotically equivalent sequences of matrices have eigenvalues that behave similarly for large $n$. 
% Indeed, denote by $\lambda_{m}(A)$ the $m$th largest eigenvalue of the matrix $A$.
% If  $\{A_n\}_{n=1}^{\infty}$ and $\{C_n\}_{n=1}^{\infty}$ asymptotically equivalent sequences of Hermitian matrices. Then for any positive integer $s$, we have 
% %Let $F(x)$ be an arbitrary continuous function. Then, we have
% \begin{equation*}
%     \lim_{n \rightarrow \infty} \frac{1}{n} \sum_{m=1}^{n}\left(\lambda^s_{m}(A_n)-\lambda^s_{m}(C_n)\right)=0. 
% \end{equation*}

% %and the eigenvalues are said to be \textit{asymptotically absolutely equally distributed} \cite{gray2006toeplitz}. In particular, for $F(r)=r$ this implies the average difference of the eigenvalues of asymptotically equivalent matrices tends to zero. 
% For more details on this we refer to \cite[Chapter 2]{gray2006toeplitz}.

% In the special case where $\{T_n\}_{n=1}^{\infty}$ is a sequence of banded Toeplitz matrices and $\{C_n\}_{n=1}^\infty$ is a sequence of circulant matrices constructed from these  Toeplitz matrices, it is known that the eigenvalues are individually close \cite{zhu2017asymptotic}. These circulant matrices are constructed by shifting and padding the rows of the banded Toeplitz matrix appropriately \cite{gray2006toeplitz}: they are circular convolutions of the generating window function $g$. 

% More precisely, suppose that the $n\times n$ banded Toeplitz matrix $T_n$ of order $n$ defined by a sequence $\{t_k, k=0, \pm1, \pm2, \hdots, \pm m\}$ where $m<<n$. Then, the corresponding circulant matrix is defined by $C_n$ whose first row 
% $\begin{pmatrix} c_0, c_1, \hdots, c_{n-1}\end{pmatrix}$ is given by 
% \begin{align}
%     c_k=\begin{cases} t_{-k}, \qquad  k=0,1, \hdots, m\\
% t_{n-k}, \qquad k=n-m, \hdots, n-1\\
% 0, \qquad {\textrm else}\end{cases}\label{eq:circ_ck}
% \end{align}

% The following result is known.



% {\color{red} analytically. explicitly compute eigenvalues of circulant matrices, instead of numerically estimating eigenvalues of finite blocks on the diagonal, or analytically estimate them with these circulant matrices. interesting on its own.

% asymptotically in a sense stated in the theorem, almost equally distributed, any continuous functions, small l1 average it goes to zero. 

% in our special cases: when toeplitz are banded that says they converge individually too. pointwise convergence. this is reference [26] they individually converge.
% }




% \begin{theorem}[again this is known, so we can just state it. \cite{gray2006toeplitz}]\label{thm:C_eigs}
% Let $C_n \in \mathbb{C}^{n\times n}$ be the circulant matrix obtained by embedding the 
% Toeplitz block $T_n$ as in \eqref{eq:circ_ck}. Then the eigenvalues of $C_n$ are given, up to permutation, by the $n$ equally spaced samples of the Laurent polynomial $t$:
% \begin{equation*}
% \operatorname{spec}(C_n)
% = 
% \bigl\{\,t(\tfrac{p}{n})\;:\,\ p=0,\dots,n-1\,\bigr\}.
% \end{equation*}
% \end{theorem}



% \subsubsection{New results for $B-$spline}
% \begin{corollary}\label{thm:C_eigs}
% Let $C^{[\ell]}_n \in \mathbb{C}^{n\times n}$ be the circulant matrix obtained by embedding the 
% Toeplitz block $T^{[\ell]}_n$ as in \eqref{eq:circ_ck}. For $n$ sufficiently large, let $\{\psi_1, \hdots, \psi_n\}$ be the ordered elements of the set \[\left\{\lambda_m (C_n^{[\ell]})\right\}_{m=0}^{n-1}=\left\{ \frac{1}{a} \sum_{r \in \mathbb{Z}}
%     \sinc^2\!\left(\frac{r - (\tfrac{p}{n} - ab\ell/2)}{a}\right)
% 	\sinc^2\!\left(\frac{r - (\tfrac{p}{n} + ab\ell/2)}{a}\right)\right\}_{p=0}^{n-1}.\]
% Then
% \begin{equation*}
%     \lim_{n \rightarrow \infty} \frac{1}{n} \sum_{m=1}^{n}\left(\lambda^s_{m}(T_n^{[\ell]})-\psi^s_{m}\right)=0. 
% \end{equation*}
    
% \end{corollary}



% The asymptotic behavior of the difference in eigenvalues between $C_n^{[\ell]}$ and $T_n^{[\ell]}$ is shown in Figure \ref{fig: eig_diff}.



% \begin{figure}[h!]
%     \centering
%          \includegraphics[width=\textwidth]{Figures/eigen_convergence.png}

%      \caption{
%      Asymptotic behavior of the difference in eigenvalues  between the Toeplitz $T_n^{[\ell]}$ and circulant $C_n^{[\ell]}$.}
%      \label{fig: eig_diff}
% \end{figure}




%\section{Frame bounds}\label{sec:bounds}
%{\color{purple}









\section{Conclusions}\label{sec:conclusion}
We conclude by connecting the results of this paper to the frame set problem for the B-spline.

We recall that if $\mathscr{G}$ is the Gram operator for a frame with frame bounds $\alpha, \beta$, then
\[\{0\}\subseteq \sigma(\mathscr{G})\subseteq\{0\}\cup [\alpha, \beta]. \]
Let  $\mathscr{G}^\dagger$ denote its Moore-Penrose pseudoinverse.
Estimates on the operator norm of the pseudo-inverse $\|\mathscr{G}^\dagger\|$ are also valuable as its reciprocal is the lower frame bound $\alpha = \|\mathscr{G}^\dagger\|^{-1}$, see \cite{adcock2019frames}. 


Furthermore, the following result can be applied to obtain frame bounds.

\begin{lemma}[Lemmas 4 and 5 in \cite{adcock2019frames}]\label{lem:adcock}
Let  $\Phi$ be a linearly independent frame with frame bounds $\alpha$ and $\beta$.
 \begin{enumerate}
 \item The truncated Gram matrix $G_n$ is invertible with $\|G^{-1}_n\|^{-1}=A_n$ and  $\|G_n\|=B_n$, where $A_n$ and $B_n$ are the frame bounds of the truncated frame $\Phi_n$.
 \item The sequences $\{A_n\}_{n\in \mathbb{Z}}$ and $\{B_n\}_{n\in \mathbb{Z}}$ are monotonically nonincreasing and monotonically nondecreasing, respectively.
 \item $B_n\leq \beta$ for all $n$ and $\lim_{n\to\infty}B_n=\beta$.
 \item $\liminf_{n}A_n>0$ iff $\Phi$ is a Riesz basis.
\end{enumerate}


\end{lemma}


\begin{remark}
\begin{enumerate}
\item Although we do not know if $\mathcal{G}(s_N,a\mathbb{Z}\times b\mathbb{Z})$ is a frame, we know that any finite section is linearly independent. This is a special case of the HRT-Conjecture which asserts that any finite collection of time-frequency shifts of a non-zero $L^2$ function is linearly independent \cite{heil1996linear, heil2007history}. Given that the HRT-conjecture is true for all compactly supported $L^2$ functions, we can prove parts 1--3 of Lemma~\ref{lem:adcock} for the Gabor systems generated by $s_N$ and all parameters $0<ab<1.$ Furthermore, we know that by the Balian-Low Theorem, $\mathcal{G}(s_N, a\mathbb{Z}\times b\mathbb{Z})$ is never a Riesz basis. As such we can conclude that $\liminf A_n=0$
\item Our results  provide upper and lower estimates on the frame bounds associated with the finite Gabor frame $\mathcal{G}(s_N,\Lambda_n)$.
\item The Bessel bound of the Gabor
system $\mathcal{G}(s_N,a\mathbb{Z}\times b\mathbb{Z})$ is the supremum of the spectrum of $T_{\infty}^{[0]}$, i.e., $B = \sup_x t^{[0]}(x)$. Suppose that $a = \tfrac{1}{p}$ with $p\in\mathbb{N}$ and $p\ge 2$. Then, Corollary \ref{corr: laurent_max_min} implies that  the Bessel bound of the Gabor system
$\mathcal{G}(g,\Lambda)$ is $\beta = p$. 
\item If $a \leq 1/2$, then the scaled and shifted $\text{sinc}^2(\frac{1}{a}(\cdot-l))$ share zeros with $\text{sinc}^2(\cdot)$ at the scaled integers $\frac{1}{a}\mathbb{Z}$. Hence, whenever $m/n$ coincides with a zero of $t^{[0]}(x)$ (for example, when $a n$ is an integer and $m$ is a multiple of $a n$), we have $\lambda_m(C_n^{[0]}) = 0$. Because the eigenvalues of the finite blocks $T_n^{[0]}$ interlace those of ${G}_n$, it must be the case that if the minimum eigenvalue of $T_n^{[0]}$ tends to zero, so does that of ${G}_n$. The minimum eigenvalue of $T_n^{[0]}$ tends to zero when $C_n^{[0]}$ has an eigenvalue of zero or tends to zero. This corresponds to when $t^{[0]}(x)$ has a zero, see Figure \ref{fig:laurent_poly}. If $a=\frac{1}{4}$ and $n=1024$, then 
\begin{equation} \label{eq: circ_eig}
    \textrm{spec}(C_n^{[0]})  = \left\{4 \sum_{r=-\infty}^{\infty} \text{sinc}^4 (\frac{p}{256}-4r )\right\}_{p=0}^{n-1}
\end{equation}
has zeros at $p=256, 512, 768$. Hence, when $a=\frac{1}{4}$ it must be that the eigenvalue and frame bound $\alpha(a,b,n)=\lambda_1(G_n)\rightarrow 0$ as $n\to \infty$.


\end{enumerate}
 
\end{remark}










It is known that for truncations of linearly independent frames, the lower frame bound must tend to zero \cite{adcock2019frames}. Thus, the following corollary can be viewed as generalizing this result to Gabor systems in general and tying them to the eigenvalues of their Gram matrices.


% \begin{lemma}[new result]
% $N>1$, let $ab<1, a<N$ and let $\Lambda =a\mathbb{Z}\times b\mathbb{Z}$. Let  $\mathcal{G}(g,\Lambda)$ be the Gabor frame associated with the $N-th$ order $B-$spline $g=s_N$ with frame bounds $A$ and $B$.
%  \begin{enumerate}
%  \item The truncated Gram matrix $G_n$ is invertible with $\|G_n\|^{-1}=A_n$ and  $\|G_n\|=B_n$, where $A_n$ and $B_n$ are the frame bounds of the truncated frame $\mathcal{G}(g,\Lambda_n)$.
%  % \item The sequences $\{A_n\}_{n\in \mathbb{Z}}$ and $\{B_n\}_{n\in \mathbb{Z}}$ are monotonically nonincreasing and monotonically nondecreasing, respectively.???? (is this true?)
%   \item $B_n\leq B$ for all $n$ and $\lim_{n\to\infty}B_n=B$.
%  \item $\liminf_{n}A_n=0$.
% \end{enumerate}
% \end{lemma}


\begin{corollary} (Lower Frame Bound and Eigenvalue Decay) \label{cor: min_eig_finite_gabor}
    If the minimum eigenvalue $\lambda_1(G_n^{[0]})$ of the of the finite block $G_n^{[0]}\in \mathbb{C}^{n\times n}$ tends to zero as $n \rightarrow\infty$, then (1) $a \in \mathbb{Q}$ and (2) the minimum eigenvalue $\lambda_1(G_n)$ of the finite Gram matrices ${G}_n\in \mathbb{C}^{n^2\times n^2}$ tends to zero as $n \rightarrow \infty$.
\end{corollary}

%Why does this imply that $a$ is rational? (because of the sinc?)

%Open question: gap between eigenvalues of $G_n$ and $\mathscr{G}$. This still evades us.




% \newpage 














In this paper, we have provided a detailed analysis of the structural properties of the Gram matrix $\mathscr{G}$ associated with Gabor systems generated by B-splines of order $N$. By exploiting the compact support of $s_N$, we established a structural hierarchy that allows for a rigorous treatment of the properties of the frame operator's finite sections.

Our primary structural result, Theorem \ref{thm:offdiagTH}, demonstrates that the blocks of the Gram matrix $G_n^{[\ell]}$ admit a Hadamard factorization $G_n^{[\ell]} = T_n^{[\ell]} \circ H_n^{[\ell]}$. This decomposition isolates the modulation-induced phase into a rank-one Hankel matrix $H_n^{[\ell]}$, while the singular values are dictated by a real-valued, symmetric, and banded Toeplitz matrix $T_n^{[\ell]}$. This result connects Gabor analysis and classical Toeplitz theory, enabling the use of Szegő-type limit theorems.




% While this work investigates the blocks of finite sections of $\mathscr{G}$ individually, the spectral relationship between the full Gram operator and its block-diagonal approximation remains an area for further study. Future research will explore:
% \begin{itemize}
%     \item The extension of the Toeplitz-dot-Hankel decomposition to more general classes of window functions in modulation spaces.
%     \item The development of fast algorithms for dual window inversion that exploit the $O(\ell^{-N})$ decay to reduce computational complexity.
%     \item The use of the circulant approximations introduced in Section \ref{sec:asymptotic} to effectively mitigate spectral pollution in the numerical computation of Gabor frames. {\color{red} this is quite a bold statement, did we indicate any strategy for this?}
% \end{itemize}

%\begin{comment}
\section{Acknowledgements}
Martin Buck and Kasso Okoudjou were partially supported by grants from the National Science Foundation under grant numbers DMS 2205771, and DMS 2309652. Christina Frederick and Alexander Stangl were supported under NSF grant DMS 2309651.

The authors also acknowledge helpful discussions with Prof. Hans Feichtinger.

The authors used AI tools to edit and polish the manuscript for spelling, grammar, and
general style. The final content and scientific conclusions remain the responsibility of the
authors
%\end{comment}
\bibliography{bib.bib}
\bibliographystyle{abbrv}




% \subsection{Removed}

% {\color{gray} this may not be needed anymore: 
% The advantage of working with circulant matrices to study the spectrum of a banded Toeplitz matrix is that their eigenvalues $\{\lambda_{m}(C_n)\}_{m=1}^n$ can be computed explicitly as the discrete Fourier Transform of a row of the $C_n$,
% \begin{equation*}
%     \lambda_{m}(C_n) = \sum_{k=0}^{n-1} c_k^{(n)} e^{- 2 \pi i m k/n}, \ m \in \{0,2, \ldots, n-1\},
% \end{equation*}
% where $c_k^{(n)}$ for $k=0,\ldots,n-1$ are entries in a row of the $C_n$. Let \[{\text{Comb}(x): = \sum_{r=- \infty }^{\infty}\delta(x-r)}\] be the \textit{Dirac comb}. We use the following relation between the DFT of a sampled function $\mathbf{f}=f(\cdot) \chi_{[-\frac{1}{2},\frac{1}{2}]}(\cdot)  \text{Comb}(\cdot)(m)$ and samples of the Fourier transform $\mathcal{F}\{f\}(m)$ for $m\in \{0,1, \ldots, n-1\}$:
% \begin{align*}
%     \text{DFT}\{\mathbf{f}\}(m) &= \big(\mathcal{F}\{f\}(
%     \cdot)*\text{sinc}( \cdot)*\text{Comb}(\cdot) \big)(\frac{m}{n}) \\
%     &= \sum_{l=- \infty}^{\infty}\big(\mathcal{F}\{f\}(\cdot)*\text{sinc}(\cdot) \big)(\frac{m}{n}-l).
% \end{align*}
% \begin{theorem} \label{thm: gii_eigs}
%     Let $C_n^{[0]} \in \mathbb{C}^{n \times n}$ be the circulant matrix constructed from the $n\times n$ Toeplitz block  $T_n^{[0]}$. Then, the eigenvalues $\{\lambda_{m}(C_n^{[0]})\}_{m=1}^n$ have the form
%     \begin{equation*}
%         \lambda_{m}(C_n^{[0]}) =\frac{1}{a} \Big(\sum_{r=-\infty}^{\infty} \textup{sinc}^2 \big(\frac{1}{a}(\frac{m}{n}-r) \big) \Big)^2=t^{[0]}(\frac{m}{n})
%     \end{equation*}
    
%     and are the square of the discrete Fourier transform of an $a$-dilated B-spline and correspond to samples of the Laurent polynomial $t^{[0]}(x)$ from Theorem \ref{thm: gii_spec}.
% \end{theorem}

% \begin{proof}
% Let 
% \begin{equation*}
%     \mathbf{g}(m) \coloneqq g (\cdot) \chi_{[-\frac{1}{2},\frac{1}{2}]} (\cdot) \text{Comb}(\cdot)(m) 
% \end{equation*}
% for $m=0, \ldots, n-1$ be the discrete set of uniform samples of $g$ on the interval $[-\frac{1}{2}, \frac{1}{2})$,
% \begin{equation*}
%     \mathbf{g} = \big(g(-\frac{1}{2}), g(-\frac{1}{2} + \frac{1}{n}), \ldots, g(-\frac{1}{2}+\frac{n-2}{n}), g(-\frac{1}{2}+\frac{n-1}{n}) \big).
% \end{equation*}
% A row of the circulant $C_n^{[0]}$ are samples of a circular convolution of $g$ with itself \cite{gray2006toeplitz}. We compute:
%     \begin{align*}
%          \lambda_m(C_n^{[0]}) &= \sum_{k=0}^{n-1}(g*g)(ak)e^{-2 \pi i km/n} \\
%         &= \sum_{k=0}^{n-1} \big(\sqrt{a}g(ax)*\sqrt{a}g(ax)\big)(k)e^{-2 \pi i km/n} \\
%         &= \big(\sum_{k=0}^{n-1}\sqrt{a}g(ak)e^{-2 \pi i km/n} \big)^2 \\
%         &= \text{DFT}\{\sqrt{a}\mathbf{g}(a \cdot)\}^2(m)
%     \end{align*}
% Thus, since $\mathcal{F}\{\sqrt{a}g(ax)\}(\cdot) = \frac{1} {\sqrt{a}}\text{sinc}^2(\frac{\cdot}{a})$, we have that
% \begin{align*}
%     \mathcal{F}\{\sqrt{a}g(a \cdot)\}(\cdot) * n\text{sinc}(n \cdot) &=
%      \frac{1}{\sqrt{a}}\text{sinc}^2(\frac{\cdot}{a})*n\text{sinc}(n \cdot) \\ &= 
%      \frac{1}{\sqrt{a}}\text{sinc}^2(\frac{\cdot}{a})
% \end{align*}
% when $n\geq \frac{1}{a}$, because
% \begin{align*}
%     \mathcal{F}\{ \frac{1}{\sqrt{a}}\text{sinc}^2(\frac{\cdot}{a})*n\text{sinc}(n \cdot)\}(\cdot) &= \sqrt{a}g(a\cdot) H_{\infty} (\frac{\cdot}{n}) \\
%     &= \sqrt{a}g(a \cdot).
% \end{align*}
% Therefore,
% \begin{align*}
%     \text{DFT}\{\sqrt{a}\mathbf{g}(a\cdot) \}(m) &= \frac{1}{\sqrt{a}} \big(\text{sinc}^2(\frac{\cdot}{a}) * \text{Comb}(\cdot) \big)(\frac{m}{n}) \\
%     &= \frac{1}{\sqrt{a}} \sum_{r=-\infty}^{\infty} \text{sinc}^2 \big(\frac{1}{a}(\frac{m}{n}-r) \big)
% \end{align*}
% The convolution with the Comb$(\cdot)$ function above is just a  periodization. The eigenvalues $\{\lambda_{m}(C_n^{[0]})\}_{m=1}^n$ corresponding to samples of the Laurent polynomial $t^{[0]}(x)$ is a standard calculation \cite{gray2006toeplitz}. It relies on the way the rows of the circulant $C_n^{[0]}$ are constructed from rows of the Toeplitz $T_n^{[0]}$:
% \begin{align*}
%     \lambda_{m}(C_n^{[0]}) &= \sum_{k=0}^{n-1}c_ke^{-2 \pi i km/n} \\
%     &= \sum_{k=0}^j g_{-k}e^{-2 \pi i mk /n} + \sum_{k=n-j}^{n-1} g_{n-k}e^{-2 \pi i m k /n} \\
%     &= \sum_{k=-j}^j g_k e^{-2 \pi i mk/n} = t^{[0]}(\frac{m}{n})
% \end{align*}
% \end{proof}
% }

% The off diagonal blocks $G_n^{[\ell]}$, $l\neq0$ are in general not Toeplitz due to the phase factor seen in the proof of Theorem \ref{thm: diag_blocks}.  Repeating the analysis of Theorem \ref{thm: gii_eigs} by relating the DFT of the vectorized or sampled function to its continuum Fourier transform leads to Theorem \ref{thm: eigs_gii'}. We also make use of the following fact: the Fourier transform of the energy-density of a function is the auto-correlation of said function.

% { {\color{red} Kasso: change this to a remark}.
% \begin{theorem}[Estimates of eigenvalues of $|G_n^{[\ell]}|^2$] \label{thm: eigs_gii'}
%     The eigenvalues $\{\lambda_{m}(C_n^{[\ell]})\}_{m=1}^n$ of the circulant matrices $C_n^{[\ell]}$ constructed from the Toeplitz blocks $|G_n^{[\ell]}|^2$ {\color{red} This is the matrix formed by takign the magnitude of each entry squared not $GG*$} are given by
%     \begin{equation*}
%         \lambda_{m}(C_n^{[\ell]}) = \sum_{k'=-\infty}^{\infty} \text{Auto-corr}\{s\}(\frac{m}{n}-k'),
%     \end{equation*}
%     where
%     \begin{equation*}
%         s(\cdot) \coloneqq \frac{1}{a}\text{sinc}^2 (\frac{\cdot}{a}) \text{sinc}^2(\frac{\cdot}{a} -bl)
%     \end{equation*}
%     and Auto-corr$\{s\}$ is the convolution of $s(\cdot)$ with itself,
%     \begin{equation*}
%         \text{Auto-corr}\{s\} = *\Big(\frac{1}{a}\text{sinc}^2 (\frac{\cdot}{a}) \text{sinc}^2(\frac{\cdot}{a} -bl)\Big) 
%     \end{equation*}
% \end{theorem}

% \begin{proof}
%     Fix an off-diagonal block consider the entries of $|G_n^{[\ell]}|^2$ for some $l\in \{0, 1, \ldots,n-1\}$. The entries have the form 
%     \begin{align*}
%         (|G_n^{[\ell]}|^2)_{0,k} &= |\int_{\mathbb{R}}g(x)g(x+ak)e^{-2 \pi i b\ell x} dx|^2 \\
%         &= |\big(\text{sinc}^2(\cdot) * (e^{2 \pi i ak (\cdot)}  \text{sinc}^2(\cdot) \big)(bl)|^2 \\
%         &= |\int_{\mathbb{R}} \text{sinc}^2(z) \text{sinc}^2(z-bl)e^{2 \pi i ak (z-bl)} dz|^2 \\
%         &= | \int_{\mathbb{R}} \text{sinc}^2(z) \text{sinc}^2(z-bl)e^{2 \pi i akz}dz|^2 \\
%         &= | \frac{1}{a} \int_{\mathbb{R}} \text{sinc}^2(\frac{z'}{a}) \text{sinc}^2(\frac{z'}{a}-bl)e^{2 \pi i kz'} dz'|^2 \\
%         &= | \mathcal{F}\{\frac{1}{a}\text{sinc}^2 (\frac{\cdot}{a}) \text{sinc}^2(\frac{\cdot}{a} -bl) \}(-k)|^2.
%     \end{align*}
%     Call $s(\cdot) \coloneqq \frac{1}{a}\text{sinc}^2 (\frac{\cdot}{a}) \text{sinc}^2(\frac{\cdot}{a} -bl)$. In the continuum where $ak$ is replaced by a $\omega \in \mathbb{R}$, we get the convolution of $\frac{1}{a}\text{sinc}^2 (\frac{\cdot}{a}) \text{sinc}^2(\frac{\cdot}{a} -bl)$ with itself,
%     \begin{align*}
%         \mathcal{F}\{|\mathcal{F}\{\frac{1}{a}\text{sinc}^2 (\frac{\cdot}{a}) \text{sinc}^2(\frac{\cdot}{a} -bl) \}|^2\} &= \text{Auto-corr}\{\frac{1}{a}\text{sinc}^2 (\frac{\cdot}{a}) \text{sinc}^2(\frac{\cdot}{a} -bl)\} \\
%         &= *\Big(\frac{1}{a}\text{sinc}^2 (\frac{\cdot}{a}) \text{sinc}^2(\frac{\cdot}{a} -bl)\Big).
%     \end{align*}
%     Thus, the eigenvalues $\lambda_{m}(C_n^{[\ell]})$ of the circulant matrix $C_n^{[\ell]}$ associated with $T_n^{[\ell]}$ are:
%     \begin{align*}
%         \lambda_{m}(C_n^{[\ell]}) &= \text{DFT}\{|\mathbf{\mathcal{F}}\{s\}|^2\}(m) \\
%         &= \mathcal{F}\{|\mathcal{F}\{s\}|^2 \}* n\text{sinc}(n \cdot) * \text{Comb}(m/n) \\
%         &= \mathcal{F}\{|\mathcal{F}\{s\}|^2 \} * \text{Comb}(m/n) \\
%         &= \sum_{k'=-\infty}^{\infty} \text{Auto-corr}\{s\}(\frac{m}{n}-k').
%     \end{align*}
% \end{proof}
% Unlike the blocks  $G_n^{[0]}$ on the diagonal, the singular values of $G_n^{[\ell]}$ are influenced by the parameter $b$ as the phase factors in the inner product  do not cancel one another.


\end{document}
